%% Compile pdflatex > biber > pdflatex > pdflatex, or: latexmk -pdf thesis
%% Set \workingdocfalse before final submission to hide 
\newif\ifworkingdoc
\workingdoctrue

\documentclass[11pt,a4paper]{article}

\usepackage[utf8]{inputenc}
\usepackage[T1]{fontenc}
\usepackage{textcomp}
\usepackage{lmodern}
\usepackage{babel}

\babelprovide[main,import]{british}

\usepackage[
  a4paper,
  top    = 2.5cm,
  bottom = 2.5cm,
  left   = 3.0cm,
  right  = 2.5cm,
]{geometry}

\usepackage{setspace}
\onehalfspacing
\setlength{\parindent}{0pt}
\setlength{\parskip}{6pt plus 2pt minus 1pt}

\usepackage{csquotes}
\usepackage[
  backend = biber,
  style   = apa,
  sorting = nyt,
  natbib  = true,
]{biblatex}
\DeclareLanguageMapping{british}{british-apa}
\addbibresource{bibliography.bib}

\usepackage{amsmath,amssymb}
\usepackage{bm}

\usepackage{booktabs}
\usepackage{longtable}
\usepackage{tabularx}
\usepackage{array}
\usepackage{multirow}
\usepackage{graphicx}
\usepackage{float}
\usepackage{rotating}
\usepackage{pdflscape}
\usepackage{enumitem}

\newcolumntype{L}[1]{>{\raggedright\arraybackslash}p{#1}}
\newcolumntype{C}[1]{>{\centering\arraybackslash}p{#1}}

\usepackage[font=small, labelfont=bf, skip=4pt]{caption}

\usepackage{xcolor}

\usepackage[
  colorlinks        = true,
  linkcolor         = blue!60!black,
  citecolor         = blue!60!black,
  urlcolor          = blue!60!black,
  pdftitle          = {Revenue Implications of the Global Minimum Tax},
  pdfauthor         = {Maurits de Haan},
  pdfsubject        = {MSc Economic Policy Thesis - Utrecht University},
  pdfkeywords       = {global minimum tax; Pillar Two; CbCR; revenue estimation; ETR; profit shifting; BEPS},
  pdfstartview      = FitH,
  bookmarksopen     = true,
  bookmarksnumbered = true,
]{hyperref}

\usepackage[colorinlistoftodos]{todonotes}

\usepackage{titlesec}
\titleformat{\section}{\large\bfseries}{\thesection}{1em}{}
\titleformat{\subsection}{\normalsize\bfseries}{\thesubsection}{1em}{}
\titleformat{\subsubsection}{\normalsize\itshape}{\thesubsubsection}{1em}{}

\makeatletter
\renewcommand{\l@section}{\@dottedtocline{1}{0em}{2.8em}}
\renewcommand{\l@subsection}{\@dottedtocline{2}{2.8em}{3.6em}}
\renewcommand{\l@subsubsection}{\@dottedtocline{3}{6.4em}{4.2em}}
\makeatother

\usepackage{fancyhdr}
\pagestyle{fancy}
\fancyhf{}
\fancyhead[L]{\small\itshape\leftmark}
\fancyhead[R]{\small\itshape de Haan (\the\year)}
\fancyfoot[C]{\small\thepage}
\renewcommand{\headrulewidth}{0.4pt}
\renewcommand{\footrulewidth}{0pt}

\fancypagestyle{plain}{%
  \fancyhf{}
  \fancyfoot[C]{\small\thepage}
  \renewcommand{\headrulewidth}{0pt}
}

%% Toggle: \workingdoctrue / \workingdocfalse (set at top of file)
%% When false, \workingnote{...} produces no output.
\newcommand{\workingnote}[1]{\ifworkingdoc\par\noindent\textcolor{red}{\textit{[Working note: #1]}}\par\fi}

\newcommand{\GMT}{GMT}
\newcommand{\GloBE}{GloBE}
\newcommand{\ETR}{ETR}
\newcommand{\SBIE}{SBIE}
\newcommand{\IIR}{IIR}
\newcommand{\UTPR}{UTPR}
\newcommand{\STTR}{STTR}
\newcommand{\CbCR}{CbCR}
\newcommand{\MNE}{MNE}
\newcommand{\pp}{\text{pp}}
\newcommand{\EUR}{\texteuro}
\newcommand{\code}[1]{\texttt{#1}}
\newcommand{\assumption}[2]{\noindent\textbf{Assumption #1.} \textit{#2}\medskip}

\begin{document}

\begin{titlepage}
  \thispagestyle{empty}
  \begin{center}

    \vspace*{1.5em}
    {\large\bfseries Utrecht University School of Economics}\\[0.3em]
    {\normalsize MSc Economic Policy}\\[3em]

    \rule{\textwidth}{1pt}\\[1.2em]

    {\LARGE\bfseries
      Decomposing Global Minimum Tax\\[0.3em]
      Revenue Uncertainty:\\[0.5em]
      \large\bfseries
      Evidence from Country-by-Country Reporting Data%
      \footnote{The copyright of this thesis rests with the author. The author is responsible for its contents and opinions expressed in the thesis. U.S.E.\ is only responsible for the academic coaching and supervision and cannot be held liable for the content. I am grateful to dr. Peter Gerbrands for valuable comments and discussions. \workingnote{check that this is complient with UU rules}}}%
    \\[1.2em]

    \rule{\textwidth}{1pt}\\[2.5em]

    {\large Maurits de Haan}\\[0.3em]
    {\normalsize\href{mailto:m.c.p.dehaan@students.uu.nl}
      {m.c.p.dehaan@students.uu.nl}}\\
    {\normalsize\href{mailto:maurits.cpdehaan@gmail.com}
      {maurits.cpdehaan@gmail.com}}\\[2.5em]

    \begin{tabular}{@{}rl@{}}
      \textbf{Supervisor:} & dr.\ Peter Gerbrands \\[0.4em]
      \textbf{Second reader:} & dr.\ Adriaan Kalwij \\[0.4em]
      \textbf{Programme:} & MSc Economic Policy \\[0.4em]
      \textbf{University:} & Utrecht University \\[0.4em]
      \textbf{Academic year:} & 2025--2026 \\[0.4em]
      \textbf{Date last updated:} & \today \\
    \end{tabular}

    \vfill

    \begin{tabular}{@{}ll@{}}
      \textbf{JEL codes:} &
        H25, H26, F23, H87 \\[0.5em]
      \textbf{Keywords:} &
        \parbox[t]{0.75\textwidth}{global minimum tax, Pillar Two,
        country-by-country reporting,\\
        revenue estimation, effective tax rate, profit shifting, BEPS}
    \end{tabular}

  \end{center}
\end{titlepage}

\newpage
\thispagestyle{plain}
\setcounter{page}{1}

\phantomsection
\addcontentsline{toc}{section}{Abstract}
\begin{center}
  {\large\bfseries Abstract}
\end{center}

\noindent

\workingnote{Need to write in last phrase}

\newpage

\tableofcontents

\newpage
\listoftables
\listoffigures

\newpage
\phantomsection
\addcontentsline{toc}{section}{List of Abbreviations}
\section*{List of Abbreviations}

\begin{tabular}{@{}ll@{}}
  \toprule
  \textbf{Abbreviation} & \textbf{Definition} \\
  \midrule
  BEA & Bureau of Economic Analysis (US) \\
  BEPS & Base Erosion and Profit Shifting \\
  CbCR & Country-by-Country Reporting \\
  CFC & Controlled Foreign Corporation \\
  CIT & Corporate Income Tax \\
  EATR & Effective Average Tax Rate \\
  EMTR & Effective Marginal Tax Rate \\
  ETR & Effective Tax Rate \\
  EU & European Union \\
  FATS & Foreign Affiliates Statistics \\
  FDI & Foreign Direct Investment \\
  GloBE & Global Anti-Base Erosion (Rules) \\
  GMT & Global Minimum Tax \\
  IIR & Income Inclusion Rule \\
  IP & Intellectual Property \\
  MNE & Multinational Enterprise \\
  NACE & Statistical Classification of Economic Activities (EU) \\
  OECD & Organisation for Economic Co-operation and Development \\
  QDMTT & Qualified Domestic Minimum Top-up Tax \\
  QRTC & Qualified Refundable Tax Credit \\
  SBIE & Substance-Based Income Exclusion \\
  SbS & Side-by-Side (package) \\
  STTR & Subject-to-Tax Rule \\
  TCJA & Tax Cuts and Jobs Act (US, 2017) \\
  UPE & Ultimate Parent Entity \\
  UTPR & Undertaxed Profits Rule \\
  \bottomrule
\end{tabular}

\newpage

\section{Introduction}\label{sec:introduction}

The Global Minimum Tax (GMT) is the most consequential reform of international corporate taxation in decades.\workingnote{will reader really understand ALP, not fiscal reader? maybe put it in a more political perspective} Agreed upon in October 2021 by 137 jurisdictions under the OECD/G20 Inclusive Framework on Base Erosion and Profit Shifting (BEPS), the GMT introduces a minimum effective tax rate (ETR) of 15\% on the profits of multinational enterprises (MNEs) with consolidated revenues exceeding \EUR 750 million. The reform was motivated by decades of intensifying profit shifting, with estimates ranging from \$310 billion to \$850 billion annually being shifted to low-tax jurisdictions \citep{torslov2023, garcia-bernardo2024, blouin2023}. Implementation of the Global Minimum Tax started in January 2024 with the European Union (EU), the United Kingdom, Japan, South Korea and several other jurisdictions. Others followed in 2025.\workingnote{show somewhere who implemented and who did not, to mosly show the status and why this research is the only thing that is possible now}

Despite the GMT's scope, its economic effects remain largely uncertain. Theoretical analyses reach different conclusions depending on the form of tax competition assumed \parencite{janeba2023}, the endogeneity of tax-haven responses \parencite{johannesen2022}, and the role of enforcement cooperation \parencite{jean2026}. Empirical revenue estimates span a wide range, from roughly \EUR100 to \EUR180~billion globally, depending on assumptions about carve-out utilisation and behavioural responses \parencite{barake2022}. The substance-based income exclusion (SBIE), which shelters returns to tangible assets and payroll from the top-up tax, may transform the minimum tax into an effective production subsidy for jurisdictions with real economic activity \parencite{schjelderup2024}.

\workingnote{GMT, tax base, torslov how muc shifted, SBIE devereux}

This thesis addresses three research questions:

\begin{enumerate}
\item How large is the GMT tax base when the volume of low-taxed profit, distribution of effective tax rates and SBIE-to-profit ratios are taken into account and how has it evolved over the 2016-2022 CbCR panel?
\item Which modelling assumptions drive the largest variation in GMT revenue estimates, and which jurisdiction characteristics predict the dominant source of uncertainty?
\item How does the SBIE effective tax floor vary across EU sectors and Member States, and does the cross-sectional pattern confirm the production-subsidy prediction of \citet{schjelderup2024}?
\end{enumerate}

This thesis makes three contributions.

First, it delivers the first systematic decomposition of GMT revenue uncertainty using the full 2016--2022 CbCR panel. Existing estimates, \citet{barake2022}, \citet{hugger2026}, and the OECD's own projections, each present a headline number without options in the modelling choices to see how certain policy and parameter choices impact the final number. By computing revenue under all combinations of five binary factors, this thesis makes the variance exactly decomposable into main effects and interactions, identifying which assumptions drive the largest uncertainty. A regression of effective floor determinants complements the decomposition by identifying which jurisdiction characteristics predict the dominant source of uncertainty.
\workingnote{look at that last sentence, since the normal regression is not what I aim to do, there is namely a deeper thought line behind this, work this out further}

Second, it provides the first empirical quantification of the SBIE's differential impact across EU sectors, computing effective tax floors and testing the \citet{schjelderup2024} prediction that the SBIE functions as a production subsidy favouring capital-intensive industries. \citet{devereux2023} compute firm-level SBIE effects but do not disaggregate by sector; this thesis fills that gap by combining CbCR aggregates with Eurostat FATS sectoral shares. The Ireland-Netherlands comparison offers a concrete test of whether the SBIE systematically advantages real-activity jurisdictions over IP-booking jurisdictions.

\workingnote{check and discuss to which extent I also include the cross validation of CbCR Tab I agains Tab IV to show the difference between implied estimate of low-taxed profits of the former agains direct measurement of the latter. Systematic comparison could namely also be an interesting comparison, since it provides a calibration estimate for the unreliable data that can be used during the GMT evaluation papers that will come from Nov 2027 onwards}

Third, it introduces a revenue pipeline framework that connects the \citet{schjelderup2024} effective floor, the \citet{hugger2026} within-jurisdiction ETR distribution, and the profit-shifting semi-elasticity literature into a single multiplicative equation. This framework makes explicit why GMT revenue estimates diverge and provides a durable structure for future research that extends beyond the specific data used here.

The thesis is organised as follows. Section~\ref{sec:literature} reviews the literature on profit shifting and the GMT. Section~\ref{sec:theory} develops the theoretical framework and hypotheses. Section~\ref{sec:methodology} presents the methodology: data, notation, the baseline revenue computation, the SBIE decomposition, the counterfactual decomposition, and the behavioural sensitivity analysis. Section~\ref{sec:results} presents the results. Section~\ref{sec:discussion} discusses the findings, implications, limitations, and concludes.

\newpage
\section{Literature review}\label{sec:literature}

\workingnote{add intro of the literature review}

\subsection{The empirical measurement of base erosion and profit shifting}\label{sec:beps}

The Global Minimum Tax was designed around the premise that MNEs shift substantial volumes of profit to low-tax jurisdictions, eroding the tax base of countries where real economic activity occurs. Quantifying the magnitude of this shifting has been a central concern of the literature for over three decades.

\subsubsection{Methodological foundations}

The foundational approach to estimating profit shifting is surveyed by \citet{heckemeyer2017}, who identify a consensus tax semi-elasticity of reported profits of approximately 0.8. This estimate, derived from a meta-regression of 27 studies, implies that a one percentage point reduction in the tax rate differential is associated with a 0.8\% increase in reported pre-tax profits. They also conclude that debt shifting accounts for roughly 30\% of observed BEPS while transfer pricing and strategic location of intellectual property account for the remaining 70\%. \citet{dharmapala2014} reaches the same conclusion based on a reveiw of 25 papers, based on the \emph{Hines-Rice approach} that regresses the logarithm of reported profits on the tax rate, controlling for real economic activity through labour and capital inputs. \citet{riedel2018} provides a subsequent review that largely confirms this consensus, noting that channel-specific studies (examining debt shifting and transfer pricing separately) tend to produce the most credible estimates.

\subsubsection{Nonlinear profit shifting}

An insight for the GMT is that profit-shifting behaviour is highly nonlinear across the tax rate distribution. \citet{dowd2017} demonstrate this using US Internal Revenue Service administrative tax return data on Controlled Foreign Corporations covering 2002--2012. Their quadratic specification reveals that the semi-elasticity at a 5\% effective tax rate is approximately 4.7, compared to only 0.7 at a 30\% ETR. This implies that a one percentage point reduction in the tax rate at the very bottom of the distribution induces nearly seven times more profit shifting than the same reduction at moderate rates. The practical consequence is: simulating an increase in all tax haven rates from 5\% to 29\% yields a \$116 billion reduction in shifted profits under the nonlinear model, roughly double the \$55--76 billion predicted by the linear model. \citet{garcia-bernardo2024} formalise this insight by developing a \emph{logarithmic} specification for the relationship between reported profits and effective tax rates. Their methodology, applied to OECD CbCR data for 2017, yields an estimate of \$850~billion in shifted profits, higher than the \$616~billion estimated by \citet{torslov2023} for 2015, partly because the log specification better captures the concentration of shifting at very low rates. Furthermore, \citet{garcia-bernardo2024} find that lower-income countries lose a larger share of total tax revenue to profit shifting than high-income countries, a finding with direct implications for the distributional assessment of the GMT.

\subsubsection{Aggregate estimates of profit shifting}

The aggregate scale of profit shifting has been estimated through two broad approaches: micro-level firm data and macro-level national accounts. The most influential macro-level study is by \citet{torslov2023}, who use Foreign Affiliates Statistics to compare the profitability of foreign-owned and locally-owned firms within each jurisdiction. Their preferred estimate is that 36\% of multinational profits (approximately \$616 billion in 2015) were shifted to tax havens, with US MNEs shifting roughly twice as much as their European counterparts. The European Union is the biggest loser of tax revenue, with domestic profits reduced by an estimated 20\% for non-haven countries. Tax haven affiliates display profit-to-wage ratios that are an order of magnitude higher than those of local firms in the same jurisdiction.

\citet{garcia-bernardo2024} produce a complementary estimate using OECD CbCR data, finding that MNEs shifted over \$850 billion in 2017. Their logarithmic specification captures the extreme concentration of shifted profits at very low ETRs and identifies the Cayman Islands, Luxembourg, Singapore, the Netherlands, Switzerland, Hong Kong, Bermuda, Puerto Rico and Ireland as the primary destinations.

\subsubsection{The double-counting critique}

A challenge to these estimates comes from \citet{blouin2023}, who demonstrate that standard data sources (US Bureau of Economic Analysis data, US IRS Form 5471, and Orbis/Bureau Van Dijk) suffer from systematic double counting and misattribution of foreign affiliate income. The equity method of accounting causes the same dollar of profit to be counted multiple times when ownership is structured through intermediate holding companies, a practice that expanded dramatically after the 1997 US ``check-the-box'' regulations. When \citet{blouin2023} replicate the revenue loss estimate of \citet{clausing2016} using the corrected measure ``Net Income Excluding Equity Income'', the estimate falls from \$77--111 billion to approximately \$11 billion, a reduction of roughly 90\%. Their reanalysis of \citet{torslov2023} suggests that the bilateral imputation of \$306 billion in shifted profits contains substantial equity income, and that removing this component reduces the total estimate by approximately half.


\subsubsection{Firm-level evidence}

\citet{johansson2017} provide large-scale firm-level evidence using the Orbis database, with 1.2 million observations across multiple countries. They find that the largest MNEs reduce their effective tax rates by 4--8.5 percentage points through a combination of profit shifting and the exploitation of tax mismatches. The estimated global corporate income tax revenue loss ranges from 4\% to 10\%. 

The representativeness of firm-level databases is itself a methodological concern. \citet{kalemli-ozcan2015} provide critical guidance on constructing nationally representative samples from Orbis. \citet{torslov2023} also shows that Orbis lacks data in low-tax jurisdictions, since publishing data is often not mandetory there. For countries like the US, Ireland and Switzerland there is no income data. They state that in 2012 only 17\% global profits can found in Orbis and that over 80\% of subsidiaries were unknown to Orbis. 

\subsection{GMT design, revenue mechanics and implementation}\label{sec:design}

\subsubsection{Revenue estimates}

\citet{barake2022} provide the first systematic estimate of GMT revenue effects under Pillar Two. Using a combination of CbCR and FATS data, they calculate that a 15\% minimum effective tax rate would generate approximately \EUR 179 billion globally\footnote{This figure assumes that all profit within a jurisdiction faces the jurisdiction-average ETR, either fully above or fully below 15\%. As \citet{hugger2026} demonstrate, 53.2\% of globally low-taxed profit lies in jurisdictions with above-15\% average ETRs, so the true GMT base is substantially larger than Barake's estimate implies. \citet{barake2022} acknowledge this limitation but lack the within-jurisdiction ETR data needed to correct it.} and \EUR 55 billion within the EU, under an IIR-first scenario. These are first-stage estimates that assume no behavioural response by MNEs. The actual revenue depends critically on whether the top-up tax is collected by the jurisdiction of the ultimate parent entity through the IIR, or by the source country through a QDMTT.

\citet{janeba2023} develop a strategic tax competition model that produces a more nuanced picture. Their finding is that the revenue effect of the GMT depends on whether governments use distortionary or lump-sum instruments to compete for real economic activity. When instruments are distortionary, the GMT reduces profit shifting but may intensify competition among non-haven countries for real investment, potentially offsetting a portion of the revenue gains.

\subsubsection{The substance-based income exclusion}

A feature of the GMT is the SBIE, which excludes from the top-up tax calculation an amount equal to a percentage of eligible payroll costs and tangible assets located in each jurisdiction. During the transition period (2024--2033), the carve-out starts at 10\% of payroll and 8\% of tangible assets, declining to 5\% of each. \citet{devereux2023} provide the first empirical assessment of the SBIE using CbCR data, finding that the average ratio of SBIE to pre-tax profit is 57\% (52\% median), declining to 41\% after the ten-year transition. This implies that the average minimum tax on GloBE income is only 6.45\% in the first transition year, rising to 8.85\% at long-run rates, well below the headline 15\%. Approximately 30\% of foreign-owned EU subsidiaries in scope have SBIE exceeding pre-tax profit, facing zero top-up tax regardless of their ETR.

\citet{schjelderup2024} analyse the economic implications of the SBIE within a formal model. They reach the conclusion that the SBIE functions as a \emph{de facto} production subsidy in low-tax jurisdictions. By tying the exclusion to tangible assets and payroll, the GMT inadvertently incentivises MNEs to increase their real economic presence in low-tax locations, thereby expanding the exclusion and reducing the top-up tax. This creates a capital demand distortion that partially offsets the reduction in transfer pricing distortions that the GMT otherwise achieves. Quantitatively, the effective reduction in the rental cost of capital (0.50--0.96\%) exceeds the reduction in the effective wage rate (up to 0.88\%), implying that the SBIE \emph{favours capital-intensive industries}.

\subsubsection{Within-jurisdiction heterogeneity}

\citet{hugger2026} provide what is the most important empirical finding for understanding the actual reach of the GMT. Using confidential ETR data at the MNE-jurisdiction level, they demonstrate that 53.2\% of all low-taxed profit (below 15\% ETR) is located in jurisdictions where the \textit{average} ETR exceeds 15\%. The within-jurisdiction dispersion in ETRs is often larger than the cross-jurisdiction dispersion. This finding implies that jurisdiction-level analyses (which use average ETRs) substantially underestimate the amount of low-taxed profit and, correspondingly, the potential revenue from the GMT.

\subsubsection{Tax incentives under Pillar Two}

\citet{ferreiraliotti2022} examine how existing tax incentives interact with the GMT. The impact varies dramatically by type: accelerated depreciation provisions have a low impact because they primarily affect timing rather than the level of taxation, whereas tax holidays and IP box regimes have a significant impact because they reduce the ETR below 15\%. The QDMTT serves as a safety valve that allows countries to retain the top-up tax domestically, but this also means that the incentive offered to the MNE is effectively nullified.

\subsection{Theoretical models of the GMT and tax competition}\label{sec:theory_lit}

\subsubsection{Welfare effects}

\citet{johannesen2022} develops the first formal model specifically designed to analyse the welfare effects of the GMT. The model features a large country (non-haven) and a small country (haven) that compete for mobile profits. The central result is that the welfare effect for the non-haven country is \textit{ambiguous}: it is positive when the profit-shifting response to taxation is strong and when the government places a low weight on the consumption of firm owners relative to public goods. The consensus semi-elasticity of 0.8 falls well below the threshold of approximately~4 required for the GMT to be unambiguously welfare-improving under Johannesen's parameterisation, though the quadratic specification of \citet{dowd2017} yields semi-elasticities above~4 for the lowest-taxed jurisdictions.

\citet{hebous2023} contribute to this debate by identifying conditions under which a minimum corporate tax can be Pareto-improving: all countries, including havens, can be made better off if the minimum tax is combined with appropriate transfers.

\subsubsection{Tax competition under the GMT}

\citet{janeba2023} model the strategic interaction between countries after GMT implementation. Their model distinguishes between profit-shifting channels (which the GMT constrains) and real-activity channels (which it does not). The GMT eliminates the incentive for havens to set rates below 15\%, but it simultaneously increases the incentive for non-havens to compete for real economic activity through subsidies, infrastructure investment, and regulatory competition.

\citet{schjelderup2024} extend this analysis by incorporating the SBIE into a two-country model. They show that the SBIE transforms the GMT from a pure anti-avoidance instrument into a hybrid policy that simultaneously discourages profit shifting (through the minimum tax) and encourages real investment in low-tax jurisdictions (through the substance exclusion). The net effect on the allocation of capital is theoretically ambiguous and depends on the elasticity of transfer pricing relative to the elasticity of real capital demand.

\subsubsection{The enforcement dilemma}

\citet{jean2026} introduce a fundamentally different concern: the GMT may undermine the international cooperation that is necessary for effective tax enforcement. In their game-theory model, two countries (high-tax and low-tax) choose both tax rates and enforcement effort. The GMT reduces the incentive for the low-tax country to cooperate on enforcement, because the top-up tax is collected by the high-tax country. When the minimum rate exceeds a critical threshold, the low-tax country prefers to withdraw from enforcement cooperation entirely. This creates a paradox: the GMT reduces profit shifting in the short run but may increase it in the long run by destroying the cooperative enforcement equilibrium.

\newpage

\section{Theoretical framework and hypotheses}\label{sec:theory}

This section develops the theoretical framework that underpins the empirical analysis. The framework has four parts. Section \ref{sec:notation} shows the notations used in this thesis. Section~\ref{sec:pipeline} introduces a revenue pipeline model that decomposes GMT revenue into three multiplicative filters, each of which reduces the gross tax base. Section~\ref{sec:floor_theory} derives the effective tax floor under Pillar Two. Section~\ref{sec:floor_sectoral} extends this to sectoral variation. Section~\ref{sec:hypotheses} states five testable hypotheses with a summary of predictions.

\subsection{Notation and variable definitions}\label{sec:notation}

Subscript $p$ denotes the parent (ultimate parent entity, UPE) jurisdiction, $j$ denotes the affiliate jurisdiction, and $t$ denotes the fiscal year ($t \in \{2016, \ldots, 2022\}$). All monetary values are in USD unless stated otherwise.

\begin{table}[H]
\centering
\caption{Variable notation}
\label{tab:notation}
\small
\begin{tabular}{@{}ll@{}}
\toprule
Symbol & Definition \\
\midrule
$\pi_{p,j,t}$ & Profit before income tax (CbCR \code{PROFIT}) \\
$\theta_{p,j,t}$ & Income tax accrued (CbCR \code{TAX\_ACCRUED}) \\
$A_{p,j,t}$ & Tangible assets excl.\ cash and equivalents (CbCR \code{ASSETS}) \\
$W_{p,j,t}$ & Total payroll / compensation of employees (CbCR \code{EARNINGS}) \\
$N_{p,j,t}$ & Number of employees, FTE (CbCR \code{EMPLOYEES}) \\
$\bar{w}_{j,t}$ & Average annual wage in jurisdiction $j$ (OECD/ILO) \\
$\tau_{p,j,t}$ & Effective tax rate $= \theta_{p,j,t} / \pi_{p,j,t}$ \\
$T$ & Global minimum tax rate $= 0.15$ \\
$\sigma_W(t)$ & SBIE payroll carve-out rate in year $t$ \\
$\sigma_A(t)$ & SBIE tangible assets carve-out rate in year $t$ \\
$\text{SBIE}_{p,j,t}$ & Substance-based income exclusion \\
$E_{p,j,t}$ & Excess profit (GloBE income minus SBIE) \\
$\delta_{p,j,t}^{\text{top}}$ & Top-up tax rate $= \max(0,\; T - \tau_{p,j,t})$ \\
$\mathbb{1}_j^{\text{Q}}$ & Indicator: $=1$ if jurisdiction $j$ has implemented a QDMTT \\
$\tau_{p,j,t}^{(k)}$ & $k$-th percentile of the ETR distribution (from Table~I) \\
$s_{j,t}^{<T}$ & Share of profit with ETR below $T$ in jurisdiction $j$ \\
$\rho_{p,j}$ & SBIE-to-profit ratio $= \text{SBIE}_{p,j} / \pi_{p,j}$ \\
\bottomrule
\end{tabular}
\end{table}

Since the CbCR data covers 2016--2022 (pre-implementation), the SBIE rates are applied counterfactually to each year's data under three rate regimes: the initial rates (10\%/8\%), the 2026 rates (7.5\%/6.2\%), and the long-run rates (5\%/5\%).
\begin{equation}\label{eq:sbie_rates_W}
\sigma_W(t) = \begin{cases} 0.10 & t = 2024 \\ 0.10 - 0.005 \cdot (t - 2024) & 2024 < t < 2033 \\ 0.05 & t \geq 2033 \end{cases}
\end{equation}
\begin{equation}\label{eq:sbie_rates_A}
\sigma_A(t) = \begin{cases} 0.08 & t = 2024 \\ 0.08 - 0.003 \cdot (t - 2024) & 2024 < t < 2033 \\ 0.05 & t \geq 2033 \end{cases}
\end{equation}

\assumption{A0}{(Population coverage). CbCR data cover MNE groups with consolidated group revenue exceeding \EUR750 million. These large MNEs account for approximately 90\% of global corporate revenue \citep{hugger2026}, but the threshold excludes smaller multinationals that may face systematically different ETRs. Since Pillar Two applies to the same \EUR750 million population, this is not a limitation for revenue estimation but matters for welfare analysis.}

\subsection{The revenue pipeline}\label{sec:pipeline}

The total GMT top-up tax revenue from jurisdiction $j$ in year $t$ can be decomposed into three sequential filters, each of which reduces the gross tax base:\footnote{The revenue pipeline decomposition is introduced in this thesis. The individual components draw on \citet{hugger2026} (Stage~1), \citet{schjelderup2024} and \citet{devereux2023} (Stage~2), and \citet{dowd2017} and \citet{garcia-bernardo2024} (Stage~3). No prior study has combined these into a single multiplicative framework that makes the sources of revenue uncertainty simultaneously visible and exactly decomposable.}

\begin{equation}\label{eq:pipeline}
R_{j,t} = \underbrace{\pi_{j,t} \cdot s_{j,t}^{<T} \cdot \bar{\delta}_{j,t}^{\text{top}}}_{\text{Stage 1: ETR filter}} \times \underbrace{\left(1 - \bar{\rho}_{j,t}\right)}_{\text{Stage 2: SBIE filter}} \times \underbrace{\left(1 - \beta_{j,t}\right)}_{\text{Stage 3: behavioural filter}}
\end{equation}

where $\pi_{j,t}$ is total MNE profit in jurisdiction $j$; $s_{j,t}^{<T}$ is the share of profit with ETR below the minimum rate $T$; $\bar{\delta}_{j,t}^{\text{top}}$ is the profit-weighted average top-up rate for the low-taxed portion; $\bar{\rho}_{j,t}$ is the SBIE-to-profit ratio applied to the low-taxed segment; and $\beta_{j,t} \in [0,1)$ is the fraction of low-taxed profit that is reallocated in response to the higher effective rate. Global revenue is $R_t = \sum_j R_{j,t}$.

\textbf{Stage~1 (ETR filter)} determines how much profit falls below the minimum rate and by how much. Two approaches exist: using the jurisdiction-average ETR (Approach~A), which treats $s_{j,t}^{<T}$ as either zero or one, or using the within-jurisdiction ETR distribution (Approach~B), which estimates $s_{j,t}^{<T}$ as a continuous share. \citet{hugger2026} show that 53.2\% of global low-taxed profit lies in jurisdictions where the average ETR exceeds 15\%, implying that Approach~A misses more than half the GMT base. The choice between these approaches is therefore the most consequential modelling decision.

\textbf{Stage~2 (SBIE filter)} determines what fraction of low-taxed profit generates a top-up tax liability. The SBIE shields a proportion $\bar{\rho}_{j,t}$ of profit from the top-up tax by excluding returns to tangible assets and payroll. \citet{devereux2023} find $\bar{\rho} = 0.57$ on average at initial carve-out rates, implying that only 43\% of GloBE income constitutes excess profit subject to the top-up tax. The SBIE-to-profit ratio varies across sectors (Section~\ref{sec:floor_sectoral}) and declines over the ten-year transition period as carve-out rates fall from 10\%/8\% to 5\%/5\%.

\textbf{Stage~3 (behavioural filter)} captures the prediction that MNEs reallocate reported profit in response to the higher effective rate. The semi-elasticity literature provides three functional forms: linear \citep{heckemeyer2017}, quadratic \citep{dowd2017}, and logarithmic \citep{garcia-bernardo2024}. All three predict that the response is concentrated at the lowest ETRs, but they disagree on its magnitude: at a 5\% ETR, the predicted profit reduction ranges from 8\% (linear) to 32\% (quadratic).

The set-up of equation~\eqref{eq:pipeline} means that revenue uncertainty compounds across stages. If Stage~1 doubles the measured base (distributional vs.\ average ETR) and Stage~2 shields 57\% of it, the combined effect is $2 \times 0.43 = 0.86$, which means that doubling the base yields only leads to a modest net increase once the SBIE is applied. The interactions between the factors have an economic meaning. If Stage~1 reveals a large volume of low-taxed profit that happens to be concentrated in capital-intensive sectors with high $\rho$, the Stage~1 $\times$ Stage~2 interaction is negative: the ETR distribution increases the base, but the SBIE shields most of the increment. Next to that, this set-up makes it possible to show which variables and scenarios matter most for revenue forecasting. If Stage~1 dominates, the priority is within-jurisdiction ETR data. If Stage~2 dominates (i.e., $\bar{\rho}$ is the main source of variation), then improving CbCR data on tangible assets and payroll is more valuable than refining ETR distributions. \workingnote{I am still doubting about this sentence, since I think it makes it more confusing than not.} 

\subsection{The effective tax floor under Pillar Two}\label{sec:floor_theory}

Under Pillar Two, the total tax on a low-taxed subsidiary ($\tau_{p,j} < T$) with positive excess profit is:
\begin{equation}\label{eq:total_tax}
\text{Tax} = \underbrace{\tau_{p,j} \cdot \pi_{p,j}}_{\text{local CIT}} + \underbrace{(T - \tau_{p,j}) \cdot (\pi_{p,j} - \text{SBIE}_{p,j})}_{\text{top-up tax}}
\end{equation}

Expanding and simplifying:
\begin{equation}\label{eq:total_tax_simplified}
\text{Tax} = T \cdot \pi_{p,j} - (T - \tau_{p,j}) \cdot \text{SBIE}_{p,j}
\end{equation}

The effective tax rate under Pillar Two is therefore:
\begin{equation}\label{eq:etr_pillar2}
\tau_{p,j}^{\text{P2}} = \frac{\text{Tax}}{\pi_{p,j}} = T - (T - \tau_{p,j}) \cdot \frac{\text{SBIE}_{p,j}}{\pi_{p,j}} = T - \delta_{p,j}^{\text{top}} \cdot \rho_{p,j}
\end{equation}

where $\rho_{p,j} \equiv \text{SBIE}_{p,j} / \pi_{p,j}$ is the SBIE-to-profit ratio.\footnote{The effective floor formula (equation~\eqref{eq:etr_pillar2}) is derived by \citet{schjelderup2024}. \citet{devereux2023} are the first to measure $\rho$ empirically, reporting an average of 57\% and a median of 52\% across EU subsidiaries in scope. The compact $\rho$ notation is introduced in this thesis to make the decomposition in Section~\ref{sec:decomposition} tractable; neither \citet{schjelderup2024} nor \citet{devereux2023} use it as an analytical device. Its analytical role in the revenue pipeline (equation~\eqref{eq:pipeline}) and factorial decomposition is original to this thesis.}

\textbf{Proposition~1.} \emph{The effective tax rate under Pillar Two is strictly below $T$ whenever $\text{SBIE}_{p,j} > 0$. The gap $T - \tau_{p,j}^{\text{P2}}$ is increasing in both the top-up rate $\delta_{p,j}^{\text{top}}$ and the SBIE-to-profit ratio $\rho_{p,j}$.} This follows directly from equation~\eqref{eq:etr_pillar2}: since $\delta_{p,j}^{\text{top}} > 0$ for any low-taxed entity and $\rho_{p,j} > 0$ whenever tangible assets or payroll are positive, the product $\delta_{p,j}^{\text{top}} \cdot \rho_{p,j}$ is strictly positive. The partial derivatives $\partial(T - \tau^{\text{P2}})/\partial \delta^{\text{top}} = \rho > 0$ and $\partial(T - \tau^{\text{P2}})/\partial \rho = \delta^{\text{top}} > 0$ confirm the monotonicity. Also, this connects to the revenue pipeline in equation~\eqref{eq:pipeline}, since the top-up tax on a single bilateral pair is equal to $\delta_{p,j}^{\text{top}} \times E_{p,j}$, where excess profit $E_{p,j} = \pi_{p,j} - \text{SBIE}_{p,j} = \pi_{p,j}(1 - \rho_{p,j})$. This is exactly Stage~2 of the pipeline: the SBIE filter transforms the gross top-up base $\delta^{\text{top}} \cdot \pi$ into the net base $\delta^{\text{top}} \cdot \pi \cdot (1 - \rho)$. The effective floor is 8.5\%, not 15\%. \citet{devereux2023} find an average SBIE-to-profit ratio of 57\% (52\% median) across EU subsidiaries in scope, with the average minimum tax on GloBE income only 6.45\% in the first transition year, rising to 8.85\% at long-run rates.

\textbf{Proposition~2.} \emph{For jurisdictions where the ETR distribution straddles $T$, the aggregate SBIE-to-profit ratio $\rho_j$ overstates the effective shielding of the GMT base. The SBIE reduces the top-up tax only on the low-taxed portion of profit, not on entities already taxed above $T$.} To see this, decompose total profit into a low-taxed segment ($\pi_j^{<T} = s_j^{<T} \cdot \pi_j$) and a segment above $T$. The SBIE is a physical quantity---it depends on tangible assets and payroll, not on the ETR. If the SBIE is allocated proportionally across segments (Assumption~A11 in Section~\ref{sec:approach_b}), the effective SBIE-to-profit ratio for the low-taxed segment equals the aggregate $\rho_j$. But if the SBIE is disproportionately concentrated in the above-$T$ segment (because high-ETR entities tend to be in capital-intensive activities with large tangible asset bases), the effective $\rho$ for the low-taxed segment is \emph{lower} than the aggregate---and the top-up tax is correspondingly \emph{higher}. Conversely, if low-taxed entities are disproportionately capital-intensive (as in Ireland's pharmaceutical sector), the effective $\rho$ exceeds the aggregate. This ambiguity motivates the sensitivity analysis in Section~\ref{sec:approach_b}, which tests the extreme cases where all SBIE is allocated to the low-taxed segment (maximum shielding) and where none is (minimum shielding).

\subsection{Sectoral variation in the effective floor}\label{sec:floor_sectoral}

The SBIE-to-profit ratio depends on the composition of economic activity. Decompose affiliates in jurisdiction $j$ into sectors $s \in \{1, \ldots, S\}$:
\begin{equation}\label{eq:rho_sectoral}
\rho_j = \frac{\text{SBIE}_j}{\pi_j} = \sum_s \omega_{j,s} \cdot \rho_{j,s}
\end{equation}
where $\omega_{j,s} = \pi_{j,s} / \pi_j$ is the profit share of sector $s$, and the sector-specific ratio is:
\begin{equation}\label{eq:rho_sector}
\rho_{j,s} = \sigma_W \cdot \frac{W_{j,s}}{\pi_{j,s}} + \sigma_A \cdot \frac{A_{j,s}}{\pi_{j,s}}
\end{equation}
The effective tax floor for sector $s$ in jurisdiction $j$ is:
\begin{equation}\label{eq:floor_sector}
\tau_{j,s}^{\text{floor}} = T - (T - \tau_j) \cdot \rho_{j,s}
\end{equation}

Capital-intensive manufacturing (high $A_{j,s}/\pi_{j,s}$) and labour-intensive services (high $W_{j,s}/\pi_{j,s}$) face lower floors than asset-light, high-profit activities (IP licensing, financial intermediation). \citet{schjelderup2024} show that the SBIE creates an asymmetric production subsidy: the effective reduction in the rental cost of capital (0.50--0.96\%) exceeds the reduction in the effective wage rate (up to 0.88\%), implying that the SBIE \emph{favours capital-intensive industries}.

Table~\ref{tab:sector_predictions} summarises the predicted direction of the SBIE effect by broad sector type. 
\begin{table}[H]
\centering
\caption{Predicted SBIE-to-profit ratios and effective floors by sector type}
\label{tab:sector_predictions}
\small
\begin{tabular}{@{}lcccc@{}}
\toprule
Sector type & $A/\pi$ & $W/\pi$ & $\rho$ prediction & Floor prediction \\
\midrule
Capital-intensive manufacturing & High & Moderate & High & Low \\
Labour-intensive services & Low & High & Moderate--High & Low--Moderate \\
IP licensing / holding & Very low & Very low & Very low & Close to $T$ \\
Financial intermediation & Moderate & Moderate & Moderate & Moderate \\
\bottomrule
\end{tabular}
\end{table}




\subsection{Summary of predictions}\label{sec:predictions}

Table~\ref{tab:predictions} collects the five hypotheses, their pipeline stages, predicted directions, and the empirical tests used in Section~\ref{sec:results}.

\begin{table}[H]
\centering
\caption{Summary of testable predictions}
\label{tab:predictions}
\small
\begin{tabular}{@{}clccp{4.5cm}@{}}
\toprule
 & Hypothesis & Stage & Predicted sign & Empirical test \\
\midrule
H1 & ETR measurement dominance & 1 & $+$ (distributional $>$ average) & Ratio of Approach~B to Approach~A revenue \\
H2 & SBIE sectoral asymmetry & 2 & Heterogeneity in $\rho_s$ ($\geq$ 3--4~\pp{} floor spread) & FATS sectoral decomposition; descriptive regression \\
H3 & Ireland--NL divergence & 2 & $\rho_{\text{IRL}} > \rho_{\text{NLD}}$, $\tau_{\text{IRL}}^{\text{floor}} < \tau_{\text{NLD}}^{\text{floor}}$ & Country-pair comparison \\
H4 & Nonlinear behavioural response & 3 & Concentrated at low ETRs; bounded aggregate & Three semi-elasticity specifications \\
H5 & Cross-validation divergence & -- & $\pi^{<T,\text{IV}} / \pi^{<T,\text{I}} > 1$ for conduits & Table~I vs.\ Table~IV ratio by jurisdiction type \\
\bottomrule
\end{tabular}
\end{table}

\newpage
\section{Methodology}\label{sec:methodology}

\subsection{Baseline revenue computation}\label{sec:baseline}

The baseline computation follows the ``tax deficit'' framework of \citet{clausing2021} and \citet{barake2022}: for each jurisdiction where MNEs face an ETR below the minimum rate, the deficit between the minimum and the actual rate is applied to the excess profit base. Two alternative ETR measurement approaches are presented side by side: (a)~jurisdiction-average ETRs and (b)~the within-jurisdiction ETR distribution from CbCR Table~I percentiles. 




\subsubsection{ETR calculation}\label{sec:etr}

The effective tax rate for each bilateral triple $(p, j, t)$ is computed from CbCR Table~I, restricting to profit-making entities:
\begin{equation}\label{eq:etr}
\tau_{p,j,t} = \frac{\theta_{p,j,t}}{\pi_{p,j,t}}, \quad \text{where } \pi_{p,j,t} > 0
\end{equation}
\assumption{A1}{(ETR numerator). Tax accrued (\code{TAX\_ACCRUED}) is used as the ETR numerator. This follows \citet{hugger2026}: accrued taxes reflect the current-year charge, unlike taxes paid (cash basis) which include timing differences. The GloBE ETR uses ``covered taxes'' which differ from both accrued and paid taxes in CbCR---it requires at least 22 adjustments from the financial accounting starting point \citep{hanlon2023}---but accrued taxes are the closest available proxy.}

\assumption{A2}{(Profit-making entities): }

For bilateral pairs where $\theta_{p,j,t} < 0$ (negative accrued taxes, approximately 5\% of observations), the ETR is set to zero. For pairs where $\tau_{p,j,t} > 1$, the ETR is winsorised at the statutory rate plus the loss-adjustment correction, following \citet{hugger2026}.

\subsubsection{SBIE calculation}\label{sec:sbie_calc}

The SBIE for each bilateral pair is:
\begin{equation}\label{eq:sbie}
\text{SBIE}_{p,j,t} = \sigma_W \cdot \hat{W}_{p,j,t} + \sigma_A \cdot A_{p,j,t}
\end{equation}
where the payroll component $\hat{W}_{p,j,t}$ is:
\begin{equation}\label{eq:payroll}
\hat{W}_{p,j,t} = \begin{cases}
W_{p,j,t} & \text{if } W_{p,j,t} \text{ is observed (\code{EARNINGS} in Table~I)} \\
N_{p,j,t} \times \bar{w}_{j,t} & \text{if } W_{p,j,t} \text{ is missing but } N_{p,j,t} \text{ is observed}
\end{cases}
\end{equation}
\assumption{A3}{(Payroll imputation): Where CbCR \code{EARNINGS} is missing, payroll is imputed as \code{EMPLOYEES} $\times$ national average wage from OECD/ILO data. This follows \citet{barake2022}. MNEs typically pay above-average wages, so this likely \emph{understates} actual payroll and therefore understates the SBIE, introducing a conservative bias (top-up tax biased upward). Luckily, \code{EARNINGS} is available for 53 of 54 parent jurisdictions in 2022 (although, this is not the case in earlier years).}

\assumption{A4}{(Tangible assets): CbCR tangible assets (\code{ASSETS}) exclude cash and cash equivalents, broadly consistent with Pillar Two eligible tangible assets. The alignment is imperfect but the discrepancy is likely small for aggregate estimates.}

\subsubsection{Top-up tax and excess profit}\label{sec:topup}

For each bilateral pair where the ETR falls below the minimum rate, the top-up tax rate is:
\begin{equation}\label{eq:topup_rate}
\delta_{p,j,t}^{\text{top}} = \max\!\left(0,\; T - \tau_{p,j,t}\right)
\end{equation}
The excess profit is GloBE income minus the SBIE, floored at zero:
\begin{equation}\label{eq:excess}
E_{p,j,t} = \max\!\left(0,\; \pi_{p,j,t} - \text{SBIE}_{p,j,t}\right)
\end{equation}
The top-up tax liability is:
\begin{equation}\label{eq:topup_tax}
\text{TopUp}_{p,j,t} = \delta_{p,j,t}^{\text{top}} \times E_{p,j,t}
\end{equation}
The top-up tax is zero whenever (a)~$\tau_{p,j,t} \geq T$, or (b)~$\text{SBIE}_{p,j,t} \geq \pi_{p,j,t}$. Condition~(b) means that even jurisdictions with ETRs below 15\% may generate no top-up tax if the SBIE is sufficiently large relative to profit. This is the mechanism through which the SBIE functions as a production subsidy \citep{schjelderup2024}. \citet{devereux2023} find that approximately 30\% of foreign-owned EU subsidiaries in scope have SBIE exceeding pre-tax profit, facing zero top-up tax regardless of their ETR.

\subsubsection{Revenue allocation}\label{sec:allocation}

The Pillar Two rule order determines which jurisdiction collects the top-up tax: (1)~QDMTT, (2)~IIR, (3)~UTPR.
\begin{align}
R_{j,t}^{\text{QDMTT}} &= \sum_p \text{TopUp}_{p,j,t} \times \mathbb{1}_j^{\text{Q}} \label{eq:qdmtt} \\
R_{p,t}^{\text{IIR}} &= \sum_j \text{TopUp}_{p,j,t} \times (1 - \mathbb{1}_j^{\text{Q}}) \label{eq:iir}
\end{align}

The net fiscal effect for jurisdiction $s$ is:
\begin{equation}\label{eq:net}
R_{s,t}^{\text{net}} = R_{s,t}^{\text{IIR}} + R_{s,t}^{\text{QDMTT}}
\end{equation}

\assumption{A5}{(QDMTT priority): Jurisdictions with a qualifying QDMTT collect the entire domestic top-up tax, taking priority over the IIR. QDMTT implementation status is sourced from the PwC Pillar Two Country Tracker. All 27 EU Member States have enacted QDMTTs via the EU Minimum Tax Directive.}

\assumption{A6}{(UTPR residual): The UTPR is treated as residual. For the EU analysis, where all Member States have implemented both IIR and QDMTT, the UTPR plays a negligible role.}

\subsubsection{Double-counting correction}\label{sec:doublecounting}

When ownership chains pass through intermediate holding companies, profit can be counted multiple times through intra-company dividends and equity method consolidation \citep{blouin2023}. The 2022 CbCR data includes an adjusted profit measure (\code{PROFIT\_ADJ}) that partially corrects for this.

\assumption{A7}{(Double-counting): }

\subsubsection{Approach A: Jurisdiction-average ETR}\label{sec:approach_a}

The first approach follows \citet{barake2022}. The jurisdiction-level aggregate ETR is the profit-weighted average across all parent jurisdictions:

\begin{equation}\label{eq:etr_agg}
\bar{\tau}_{j,t} = \frac{\sum_p \theta_{p,j,t}}{\sum_p \pi_{p,j,t}}
\end{equation}

Under this approach, the entire profit of jurisdiction $j$ is treated as either above or below the GMT threshold. \citet{hugger2026} show that 53.2\% of global low-taxed profit lies in jurisdictions where the average ETR exceeds 15\%, implying that Approach~A misses more than half the GMT base.

\subsubsection{Approach B: Within-jurisdiction ETR distribution}\label{sec:approach_b}

The second approach 

\paragraph{Step 1: Aggregate bilateral percentiles to jurisdiction level.}
CbCR Table~I reports percentiles (P5, P25, P50, P75, P95) of the ETR distribution for each bilateral pair. The profit-weighted jurisdiction-level percentile is:

\begin{equation}\label{eq:jurisdiction_percentile}
\tau_{j,t}^{(k)} = \frac{\sum_p \pi_{p,j,t} \cdot \tau_{p,j,t}^{(k)}}{\sum_p \pi_{p,j,t}}, \quad k \in \{5, 25, 50, 75, 95\}
\end{equation}

\assumption{A8}{(Parent-independence): The ETR distribution in affiliate jurisdiction $j$ does not vary systematically by parent jurisdiction $p$. This follows \citet{hugger2026}. It is violated when bilateral tax treaties create parent-specific withholding rates, but these are typically small relative to the CIT component.}

\paragraph{Step 2: Fit a piecewise-uniform distribution.}
Following \citet{hugger2026}, 

\begin{equation}\label{eq:piecewise}
f_{j,t}(\tau) = \begin{cases}
\text{point mass at } \tau_{j,t}^{(5)} & \text{for } \tau < \tau_{j,t}^{(5)} \\[4pt]
\frac{0.20}{\tau_{j,t}^{(25)} - \tau_{j,t}^{(5)}} & \text{for } \tau_{j,t}^{(5)} \leq \tau < \tau_{j,t}^{(25)} \\[4pt]
\frac{0.25}{\tau_{j,t}^{(50)} - \tau_{j,t}^{(25)}} & \text{for } \tau_{j,t}^{(25)} \leq \tau < \tau_{j,t}^{(50)} \\[4pt]
\frac{0.25}{\tau_{j,t}^{(75)} - \tau_{j,t}^{(50)}} & \text{for } \tau_{j,t}^{(50)} \leq \tau < \tau_{j,t}^{(75)} \\[4pt]
\frac{0.20}{\tau_{j,t}^{(95)} - \tau_{j,t}^{(75)}} & \text{for } \tau_{j,t}^{(75)} \leq \tau < \tau_{j,t}^{(95)} \\[4pt]
\text{point mass at } \tau_{j,t}^{(95)} & \text{for } \tau \geq \tau_{j,t}^{(95)}
\end{cases}
\end{equation}

\assumption{A9}{(Piecewise-uniform): . This is the central distributional assumption, following \citet{hugger2026}.  The assumption is most strained for jurisdictions with bimodal ETR distributions---\citet{hugger2026} find that within-jurisdiction ETR dispersion (P5--P95) exceeds 20 percentage points for more than half of all jurisdictions. The asymmetric shift in Step~3 (equation~\eqref{eq:shift}) partially preserves this bimodality for jurisdictions identified in the OECD Investment Tax Incentives Database, but the piecewise-uniform assumption remains a smoothing approximation.}

\paragraph{Step 3: Shift the distribution to match the profit-weighted average.}

\begin{equation}\label{eq:shift}
\tilde{\tau}_{j,t}^{(k)} = \tau_{j,t}^{(k)} + \left(\bar{\tau}_{j,t} - \bar{\tau}_{j,t}^{\text{dist}}\right)
\end{equation}

where $\bar{\tau}_{j,t}^{\text{dist}}$ is the mean of the piecewise-uniform distribution from Step~2. The shift is bounded: $\tilde{\tau}_{j,t}^{(k)} \geq 0$. Following \citet{hugger2026}, the shift is applied asymmetrically for jurisdictions with broad zero-tax incentives.

\paragraph{Computing the share of low-taxed profit.}
The share of profit with an ETR below $T$ is:

\begin{equation}\label{eq:lowtax_share}
s_{j,t}^{<T} = \int_0^{T} \tilde{f}_{j,t}(\tau) \, d\tau
\end{equation}

Under the piecewise-uniform assumption, this reduces to a sum of rectangular areas. If $T$ falls between $\tilde{\tau}_{j,t}^{(k)}$ and $\tilde{\tau}_{j,t}^{(k+1)}$:

\begin{equation}\label{eq:lowtax_share_pw}
s_{j,t}^{<T} = \underbrace{\sum_{m < k} (\text{mass in segment } m)}_{\text{segments entirely below } T} + \underbrace{\frac{T - \tilde{\tau}_{j,t}^{(k)}}{\tilde{\tau}_{j,t}^{(k+1)} - \tilde{\tau}_{j,t}^{(k)}} \times (\text{mass in segment } k)}_{\text{partial segment}}
\end{equation}

The volume of low-taxed profit is:

\begin{equation}\label{eq:lowtax_profit}
\pi_{j,t}^{<T} = s_{j,t}^{<T} \times \sum_p \pi_{p,j,t}
\end{equation}

\assumption{A10}{(Proportional profit): Profit is distributed in proportion to the mass of MNE subgroups across the ETR distribution. In reality, subgroups with lower ETRs tend to be larger (concentrated profits in tax-haven affiliates), so this likely \emph{understates} the volume of low-taxed profit. The profit-weighted percentiles in Step~1 partially correct for this but do not fully resolve the within-segment weighting problem.}

\paragraph{Revenue recomputation under the distribution.}
Under Approach~B, the top-up tax becomes:

\begin{equation}\label{eq:topup_distrib}
\text{TopUp}_{j,t}^{\text{B}} = \sum_p \left[\bar{\delta}_{j,t}^{\text{top}} \times \max\!\left(0,\; s_{j,t}^{<T} \cdot \pi_{p,j,t} - \text{SBIE}_{p,j,t}^{<T}\right)\right]
\end{equation}

where $\bar{\delta}_{j,t}^{\text{top}}$ is the profit-weighted average top-up rate for the low-taxed portion:

\begin{equation}\label{eq:avg_topup_rate}
\bar{\delta}_{j,t}^{\text{top}} = T - \frac{\int_0^{T} \tau \cdot \tilde{f}_{j,t}(\tau) \, d\tau}{\int_0^{T} \tilde{f}_{j,t}(\tau) \, d\tau}
\end{equation}

and the SBIE allocated to the low-taxed portion is:

\begin{equation}\label{eq:sbie_lowtax}
\text{SBIE}_{p,j,t}^{<T} = s_{j,t}^{<T} \times \text{SBIE}_{p,j,t}
\end{equation}

\assumption{A11}{(Proportional SBIE): The SBIE is distributed proportionally to the low-taxed profit share. In reality, the SBIE may be disproportionately concentrated in either segment. Sensitivity analysis tests the extreme cases where all SBIE is allocated to the low-taxed segment (upper bound on shielding) and where none is (lower bound).}


\subsection{SBIE decomposition and the effective tax floor}\label{sec:sbie_decomp}


CbCR does not report by sector. The sectoral decomposition requires Eurostat FATS for the within-jurisdiction allocation of payroll and tangible assets across sectors.

For each EU Member State $j$ and sector $s$, the sectoral SBIE is:

\begin{equation}\label{eq:sbie_sectoral}
\widehat{\text{SBIE}}_{j,s} = \sigma_W \cdot \left(\frac{\text{PersC}_{j,s}^{\text{FATS}}}{\sum_{s'} \text{PersC}_{j,s'}^{\text{FATS}}}\right) \cdot W_j^{\text{CbCR}} + \sigma_A \cdot \left(\frac{\hat{K}_{j,s}}{\sum_{s'} \hat{K}_{j,s'}}\right) \cdot A_j^{\text{CbCR}}
\end{equation}

FATS provides the \emph{sectoral shares}; CbCR provides the \emph{aggregate levels}.

\assumption{A13}{(FATS--CbCR proxy). The sectoral shares of employment and personnel costs in Eurostat FATS (for foreign-controlled enterprises $>$50\% foreign-owned) are a valid proxy for the sectoral shares of CbCR \code{EARNINGS} and \code{EMPLOYEES} (for MNE groups $>$\EUR750M revenue). The CbCR population is a subset of the FATS population.}

\assumption{A13a}{(FATS series break). Eurostat inward FATS underwent a structural break in 2021. The pre-2021 dataset (\code{fats\_g1b\_08}) uses 23 legacy indicator codes, 255 control-country categories, and a broad activity aggregate. The post-2021 dataset (\code{fats\_activ}) uses 19 revised indicator codes, 50 control-country categories, and a different aggregate. Key concept changes include ``turnover'' becoming ``net turnover'' and ``personnel costs'' becoming ``employee benefits expense.'' The sectoral analysis uses the post-2021 structure (2021--2023) for all cross-sectional results. Pre-2021 data is used only for trend context and is not pooled with post-2021 data without explicit harmonisation.}

\assumption{A14}{(Tangible asset imputation). The sectoral share of tangible assets is proportional to value added or employment, scaled by sector-level capital intensity ratios from the OECD STAN database. FATS does not report tangible assets directly. This is the weakest link in the sectoral decomposition and should be interpreted as indicative.}

The effective tax floor evolves over the SBIE transition period (2024--2033):

\begin{equation}\label{eq:floor_transition}
\tau_{j,s}^{\text{floor}}(t) = T - (T - \tau_j) \cdot \left[\sigma_W(t) \cdot \frac{W_{j,s}}{\pi_{j,s}} + \sigma_A(t) \cdot \frac{A_{j,s}}{\pi_{j,s}}\right]
\end{equation}

The output is a matrix of effective tax floors: 27 EU Member States $\times$ 15--20 NACE sectors $\times$ 10 transition years.





%% ═══════════════════════════════════════════════════════════════════════
%%  5. RESULTS
%% ═══════════════════════════════════════════════════════════════════════


[To be completed after data analysis.]

\subsection{Descriptive statistics}\label{sec:descriptive}

[CbCR coverage, ETR distributions by jurisdiction, SBIE-to-profit ratios.]

\subsection{Baseline revenue estimates}\label{sec:baseline_results}

[Approach A vs.\ Approach B comparison. EU-27 revenue table by Member State.]

\subsection{Cross-validation results}\label{sec:crossval_results}

[Table I vs.\ Table IV comparison. Validation ratios by jurisdiction.]

\subsection{SBIE sectoral decomposition}\label{sec:sbie_results}

[Effective tax floors by sector and Member State. Ireland--Netherlands comparison. Transition path.]

\subsection{Counterfactual decomposition}\label{sec:decomp_results}

[Marginal contributions. Factor interactions. Revenue uncertainty map.]

\subsection{Behavioural sensitivity}\label{sec:behav_results}

[Revenue estimates under three semi-elasticity specifications. Scenario matrix.]

\subsection{Robustness checks}\label{sec:robustness}

[Alternative distributional assumptions. Double-counting sensitivity. SBIE allocation alternatives.]


%% ═══════════════════════════════════════════════════════════════════════
%%  6. DISCUSSION AND CONCLUSION
%% ═══════════════════════════════════════════════════════════════════════
\newpage
\section{Discussion and conclusion}\label{sec:discussion}

\subsection{Discussion}\label{sec:discussion_main}

[Position findings relative to \citet{barake2022}, \citet{hugger2026}, and OECD projections. Relate to theoretical predictions of \citet{schjelderup2024} and \citet{johannesen2022}.]

\subsection{Policy implications}\label{sec:implications}

[Implications for policymakers: which factor drives the largest uncertainty. Implications for the SBIE design debate. Implications for QDMTT vs.\ IIR allocation.]

\subsection{Limitations}\label{sec:limitations}

The methodology captures first-order revenue effects through a microsimulation framework. Several dimensions that the literature identifies as potentially consequential lie outside its scope.

\textbf{CbCR-to-GloBE approximation.} The CbCR ETR is an approximation of the GloBE ETR. CbCR taxes accrued do not include the deferred tax adjustment, and CbCR does not distinguish QRTCs from non-qualified credits. The direction and magnitude of the approximation error vary by jurisdiction (Assumptions~A1, A4a).

\textbf{Real economic responses.} The semi-elasticity specifications model reallocation of \emph{reported} profits (paper shifting), but the GMT may also trigger reallocation of real economic activity. \citet{suarezserrato2018} find that the repeal of US Section~936 reduced global investment by 10\% and US employment by 6.7\%. \citet{bilicka2022} find that the UK Worldwide Debt Cap reduced total assets by 30.6\% among treated foreign MNEs. If the GMT triggers analogous responses, the SBIE itself would shrink, amplifying the top-up tax in a second-round effect not captured here (Assumption~A17).

\textbf{Government strategic responses.} \citet{janeba2023} show that the GMT may intensify competition for real economic activity through subsidies and regulatory competition. The introduction of QDMTTs is itself a strategic response that fundamentally changes revenue allocation.

\textbf{Enforcement dynamics.} \citet{jean2026} model the interaction between the GMT and endogenous enforcement effort, finding that the GMT may undermine international enforcement cooperation when the minimum rate exceeds a critical threshold.

\textbf{Static computation.} The methodology applies a static year-by-year computation without inter-temporal loss carry-forward or dynamic adjustment. GloBE rules allow net losses to reduce GloBE income in subsequent years (Assumption~A4c).

\textbf{No Pillar One interaction.} The analysis does not model Pillar One (Amount~A), which would reallocate a portion of residual profits to market jurisdictions.

\textbf{Welfare analysis.} Revenue estimates do not map directly to welfare. \citet{johannesen2022} demonstrates theoretical ambiguity that depends on the profit-shifting semi-elasticity and redistributive preferences.

\subsection{Conclusion}\label{sec:conclusion}

[Formulate answer to main research questions. Summarise the main argument and its theoretical and practical relevance. Discuss validity and reliability of findings. Nuance conclusions. Suggestions for future research.]


%% ═══════════════════════════════════════════════════════════════════════
%%  REFERENCES
%% ═══════════════════════════════════════════════════════════════════════
\newpage
\printbibliography[heading=bibintoc]


%% ═══════════════════════════════════════════════════════════════════════
%%  APPENDICES
%% ═══════════════════════════════════════════════════════════════════════
\newpage
\appendix
\renewcommand{\thesection}{\Alph{section}}
\renewcommand{\thetable}{\Alph{section}.\arabic{table}}
\renewcommand{\thefigure}{\Alph{section}.\arabic{figure}}
\renewcommand{\theequation}{\Alph{section}.\arabic{equation}}
\setcounter{table}{0}
\setcounter{figure}{0}
\setcounter{equation}{0}

%% ─────────────────────────────────────────────────────────────────────
%%  APPENDIX A: Summary of assumptions
%% ─────────────────────────────────────────────────────────────────────
\section{Summary of assumptions}\label{app:assumptions}
\setcounter{table}{0}

\begin{longtable}{@{}lp{10cm}l@{}}
\caption{Summary of assumptions}\label{tab:assumptions}\\
\toprule
Code & Assumption & Section \\
\midrule
\endfirsthead
\multicolumn{3}{c}{\textit{Table~\ref{tab:assumptions} continued}}\\
\toprule
Code & Assumption & Section \\
\midrule
\endhead
\bottomrule
\endlastfoot
A0 & CbCR population (\EUR750M threshold) covers $\sim$90\% of global corporate revenue & \ref{sec:notation} \\
A1 & Tax accrued is the appropriate ETR numerator; CbCR ETRs lower than GloBE ETRs for accelerated-depreciation jurisdictions, higher for QRTC-intensive ones & \ref{sec:etr} \\
A2 & Only profit-making entities (panel~AI) enter the ETR calculation & \ref{sec:etr} \\
A3 & Missing payroll imputed as employees $\times$ national average wage & \ref{sec:sbie_calc} \\
A4 & CbCR tangible assets approximate Pillar Two eligible tangible assets & \ref{sec:sbie_calc} \\
A5 & QDMTT jurisdictions collect the entire domestic top-up tax & \ref{sec:allocation} \\
A6 & UTPR is treated as residual and not modelled separately for EU analysis & \ref{sec:allocation} \\
A7 & Double-counting correction uses \code{PROFIT\_ADJ} where available; median correction 32.4\% & \ref{sec:doublecounting} \\
A8 & ETR distribution does not vary by parent jurisdiction within an affiliate & \ref{sec:approach_b} \\
A9 & ETRs are piecewise-uniformly distributed between percentiles; bimodal distributions partially preserved via asymmetric shift & \ref{sec:approach_b} \\
A10 & Profit distributed proportionally to MNE subgroup mass across ETR distribution; likely understates low-taxed profit & \ref{sec:approach_b} \\
A11 & SBIE distributed proportionally to low-taxed profit share; tested via extreme-case sensitivity & \ref{sec:approach_b} \\
A12 & Tables I and IV are consistent measurements of the same population; CHE excluded & \ref{sec:crossval} \\
A13 & FATS sectoral employment/payroll shares proxy CbCR sectoral composition & \ref{sec:sbie_decomp} \\
A13a & FATS series break at 2021; post-2021 structure used for cross-sectional analysis & \ref{sec:sbie_decomp} \\
A14 & Tangible asset shares imputed from external capital intensity ratios (weakest link) & \ref{sec:sbie_decomp} \\
A15 & Imported semi-elasticities are structurally stable; supported by \citeauthor{beer2020} and TCJA evidence & \ref{sec:behavioural} \\
A16 & Garcia-Bernardo and Dowd parameters taken as given without re-estimation & \ref{sec:behavioural} \\
A16a & Dowd statutory-rate coefficients applied to CbCR effective rates; Dowd estimates are upper bounds & \ref{sec:behavioural} \\
A17 & SBIE computed on pre-response values; no second-round real adjustment & \ref{sec:behavioural} \\
A18 & No MNE restructuring to avoid GMT (sub-threshold splitting, safe harbours) & \ref{sec:behavioural} \\
\end{longtable}


%% ─────────────────────────────────────────────────────────────────────
%%  APPENDIX B: Variable concordance
%% ─────────────────────────────────────────────────────────────────────
\newpage
\section{Variable concordance}\label{app:concordance}
\setcounter{table}{0}

\begin{landscape}
\begin{longtable}{@{}p{2.5cm}p{3cm}p{2cm}p{2.5cm}p{3cm}p{3cm}@{}}
\caption{Variable concordance: model to data}\label{tab:concordance_table}\\
\toprule
Model variable & Data source & Table/File & Dimension filter & Measure code & Notes \\
\midrule
\endfirsthead
\multicolumn{6}{c}{\textit{Table~\ref{tab:concordance_table} continued}}\\
\toprule
Model variable & Data source & Table/File & Dimension filter & Measure code & Notes \\
\midrule
\endhead
\midrule
\multicolumn{6}{r}{\textit{Continued on next page}}\\
\endfoot
\bottomrule
\endlastfoot

$\pi_{p,j,t}$ & OECD CbCR & Table I & \code{PANELAI}, \code{\_Z} (ETR), \code{\_Z} (stat.\ op.) & \code{PROFIT} & Profit-makers only; USD \\
$\pi_{p,j,t}^{\text{adj}}$ & OECD CbCR & Table I & Same as above & \code{PROFIT\_ADJ} & Adjusted for double-counting \\
$\theta_{p,j,t}$ & OECD CbCR & Table I & Same as above & \code{TAX\_ACCRUED} & Current-year charge \\
$A_{p,j,t}$ & OECD CbCR & Table I & Same as above & \code{ASSETS} & Tangible assets excl.\ cash \\
$W_{p,j,t}$ & OECD CbCR & Table I & Same as above & \code{EARNINGS} & Total payroll \\
$N_{p,j,t}$ & OECD CbCR & Table I & Same as above & \code{EMPLOYEES} & FTE count \\
$\bar{w}_{j,t}$ & OECD wages / ILO & CSV files & Latest available year & -- & For payroll imputation \\
$\tau_{p,j,t}^{(k)}$ & OECD CbCR & Table I & \code{PANELAI}, \code{P5}/\code{P25}/\code{P50}/\code{P75}/\code{P95} & \code{ETR} & Bilateral ETR percentiles \\
$\pi_{j,t}^{<T}$ (direct) & OECD CbCR & Table IV & \code{RATE\_0-15}, \code{TRSUB} & \code{PROFIT} & Direct GMT-base measurement \\
$A_{j,t}^{<T}$ (direct) & OECD CbCR & Table IV & \code{RATE\_0-15}, \code{TRSUB} & \code{ASSETS} & SBIE denominator for low-taxed segment \\
$W_{j,t}^{<T}$ (direct) & OECD CbCR & Table IV & \code{RATE\_0-15}, \code{TRSUB} & \code{EARNINGS} & SBIE denominator for low-taxed segment \\
$\text{PersC}_{j,s}^{\text{FATS}}$ & Eurostat FATS & \code{fats\_activ} & By \code{nace\_r2}, \code{c\_ctrl} & Personnel costs & Sectoral payroll proxy \\
$\mathbb{1}_j^{\text{Q}}$ & PwC tracker & CSV & QDMTT status & -- & Binary \\
$\tau_{j}^{\text{STR}}$ & OECD CIT rates & CSV & Combined central + sub-central & \code{CIT\_RATE} & For ETR capping \\

\end{longtable}
\end{landscape}


%% ─────────────────────────────────────────────────────────────────────
%%  APPENDIX C: CbCR table capabilities
%% ─────────────────────────────────────────────────────────────────────
\newpage
\section{CbCR table capabilities}\label{app:cbcr_tables}
\setcounter{table}{0}

\begin{table}[H]
\centering
\caption{CbCR table capabilities for the methodology}
\label{tab:capabilities}
\small
\begin{tabular}{@{}lp{5.5cm}p{5.5cm}@{}}
\toprule
Table & Can do & Cannot do \\
\midrule
I & Bilateral financial aggregates for 58$\times$240 pairs; ETR percentiles (P5--P95); full 2016--2022 panel; distinguishes profit-makers & No ETR bracket decomposition; no sectoral breakdown; percentiles may be suppressed for small cells \\
II & Interquartile means (robust to outliers); 46 parent jurisdictions & No aggregates (only central tendency); fewer parents than Table~I; no distributional shape \\
III & Entity counts by ETR bracket at MNE-group level; activity shares & No financial aggregates; MNE-group ETR misses within-group heterogeneity \\
IV & \textbf{Direct measurement of profit in 0--15\% ETR bracket at subsidiary level}; all financial variables by bracket & Fewer jurisdictions (49$\times$154); Switzerland anomaly; no percentile distributions within brackets \\
V & MNE size distribution (percentiles of revenue, assets, employees) & No bilateral breakdown; 40 parents only \\
VI & Financial aggregates by MNE size class; 2022 only & Only one year; 39 parents; no time series \\
\bottomrule
\end{tabular}
\end{table}


%% ─────────────────────────────────────────────────────────────────────
%%  APPENDIX D: Scenario matrix
%% ─────────────────────────────────────────────────────────────────────
\newpage
\section{Scenario matrix}\label{app:scenarios}
\setcounter{table}{0}

\begin{table}[H]
\centering
\caption{Revenue estimates under implementation regime $\times$ behavioural response}
\label{tab:scenario_matrix}
\small
\begin{tabular}{@{}lccc@{}}
\toprule
& \multicolumn{3}{c}{Behavioural response} \\
\cmidrule(l){2-4}
Implementation regime & None (mechanical) & Linear ($\varepsilon = 0.8$) & Quadratic/Log \\
\midrule
Full global Pillar Two & $R_{11}$ & $R_{12}$ & $R_{13}$ \\
Status quo (incl.\ SbS) & $R_{21}$ & $R_{22}$ & $R_{23}$ \\
EU-only implementation & $R_{31}$ & $R_{32}$ & $R_{33}$ \\
\bottomrule
\end{tabular}
\end{table}


%% ─────────────────────────────────────────────────────────────────────
%%  APPENDIX E: Additional tables and figures
%% ─────────────────────────────────────────────────────────────────────
\newpage
\section{Additional tables and figures}\label{app:additional}

[Robustness tables, supplementary figures, detailed Member State results.]


%% ─────────────────────────────────────────────────────────────────────
%%  APPENDIX F: Data sources
%% ─────────────────────────────────────────────────────────────────────
\newpage
\section{Data sources}\label{app:data}
\setcounter{table}{0}

This appendix documents all data sources used in the analysis, their structure, coverage, and limitations.

\subsection{OECD CbCR data}\label{app:cbcr}

Under the OECD/G20 BEPS Action~13, large MNE groups with consolidated revenue exceeding EUR~750~million are required to file Country-by-Country Reports detailing their economic activity and tax payments in each jurisdiction. The OECD aggregates and anonymises these reports into jurisdiction-level statistics published in the \emph{OECD Corporate Tax Statistics} database.\footnote{Data available at \url{https://www.oecd.org/en/topics/sub-issues/corporate-tax-statistics.html}. The individual table dataflows follow the SDMX format \code{OECD.CTP.TPS,DSD\_CBCR@DF\_CBCR[I--VI],1.1} and can be accessed programmatically via the OECD SDMX API.} The tables cover fiscal years 2016--2022 and contain approximately 2.47 million rows (330~MB uncompressed CSV).

\subsubsection{Common column structure}

All six CbCR files share fifteen columns: \code{DATAFLOW} (SDMX dataflow identifier), \code{REF\_AREA} (parent jurisdiction, ISO~3166-1 alpha-3), \code{COUNTERPART\_AREA} (affiliate jurisdiction), \code{FREQ} (always annual), \code{MEASURE}, \code{UNIT\_MEASURE}, \code{STATISTICAL\_OPERATION}, \code{AGGREGATION\_TYPE}, \code{PROFIT\_GROUPING}, \code{EFFECTIVE\_TAX\_RATE}, \code{MNE\_SIZE}, \code{TIME\_PERIOD}, \code{OBS\_VALUE} (in USD), \code{DECIMALS}, and \code{OBS\_STATUS}.

\subsubsection{Key measures}

Financial aggregates: \code{PROFIT} (profit before income tax), \code{PROFIT\_ADJ} (adjusted for double-counting, 2021--2022 only), \code{TAX\_ACCRUED} (current-year charge), \code{TAX\_PAID} (cash basis), \code{TOT\_REV} (total revenue), \code{UPR} (unrelated-party revenue), \code{RPR} (related-party revenue), \code{ASSETS} (tangible assets excl.\ cash), \code{EARNINGS} (total payroll), \code{STATED\_CAPITAL}. Count variables: \code{EMPLOYEES} (FTE), \code{ENTITIES\_COUNT}, \code{CBCR\_COUNT} (MNE groups). Derived: \code{ETR} (effective tax rate, Table~I only). Business-activity counts: \code{ACT\_ADMIN}, \code{ACT\_HOLDING}, \code{ACT\_IP}, \code{ACT\_MAN}, \code{ACT\_SALES}, and others.

\subsubsection{ETR bracket codes}

The ETR brackets in Tables~III and~IV include: \code{RATE\_0-15} (the GMT-critical bracket: all profit in this range triggers the Pillar Two top-up tax), \code{RATE\_0-10} (most aggressive avoidance), \code{RATE\_10} (10--15\%), \code{RATE\_15} (15--20\%), \code{RATE\_20-30}, \code{RATE\_30} (above 30\%), \code{NEGATIVE\_PROFIT} (loss-makers), and \code{NEGATIVE\_TAX\_ACCRUED} (positive profit but negative tax).

\subsubsection{Cross-table summary}

\begin{table}[H]
\centering
\caption{Cross-table summary of OECD CbCR Tables I--VI}
\label{tab:app_crosssummary}
\footnotesize
\begin{tabular}{@{}lrrrrrr@{}}
\toprule
 & \textbf{I} & \textbf{II} & \textbf{III} & \textbf{IV} & \textbf{V} & \textbf{VI} \\
 & Agg.\ Tot. & IQ Means & ETR (MNE) & ETR (Sub.) & Size Dist. & Size Agg. \\
\midrule
Rows & 1,124k & 437k & 699k & 101k & 4k & 106k \\
Years & 2016--22 & 2017--22 & 2017--22 & 2016--22 & 2016--22 & 2022 \\
Parents & 58 & 46 & 46 & 49 & 40 & 39 \\
Affiliates & 240 & 238 & 238 & 154 & --- & 215 \\
Agg.\ type & TOTAL & TOTAL & TRMNE & TRSUB & TRMNE & MNESIZE \\
ETR brackets & No & No & Yes (12) & Yes (11) & No & No \\
\bottomrule
\end{tabular}
\end{table}

\subsubsection{Year-by-year coverage}

The most recent release (fiscal year 2022) reports data from 54 parent jurisdictions, representing approximately 98\% of all CbCR filings \citep{oecd2025corptaxstats}. Coverage grew from 26 parent jurisdictions in 2016 to 54 in 2022. The 2021 vintage shows a spike in missing values (30,418 NaN rows in Table~I), reflecting the transition to the new SDMX format.

\subsubsection{Switzerland anomaly in Table~IV}

Switzerland has anomalous values in Table~IV for fiscal years 2017, 2020, and 2022: reported profit values approximately 8,000--10,000 times larger than the corresponding Table~I totals (e.g., USD~944~trillion vs.\ USD~109~billion in 2022). The years 2018 and 2021 are consistent. This reflects an error in the Swiss national CbCR submission. \textbf{Switzerland is excluded from Table~IV ETR bracket computations for 2017, 2020, and 2022.}

\subsection{OECD statutory corporate income tax rates}\label{app:cit}

\textbf{File:} \code{raw/oecd\_tax\_rates/oecd\_statutory\_cit\_rates.csv}\\
\textbf{Rows:} 13,273 \quad \textbf{Years:} 2000--2025 \quad \textbf{Countries:} 144

Key measures: \code{CIT} (combined central + sub-central rate), \code{CIT\_C} (central only), \code{CIT\_SCG} (sub-central). Targeting dimension: \code{ST} (standard), \code{RED} (reduced), \code{INC} (increased).

\subsection{OECD forward-looking effective tax rates}\label{app:etr}

Three files provide forward-looking EATR and EMTR measures:

\begin{itemize}[noitemsep]
\item \textbf{Baseline} (no incentives): 29,489 rows, 103 countries, 2017--2024
\item \textbf{Income-based incentives} (IP boxes): 5,784 rows, 55 countries, 2000--2024
\item \textbf{Expenditure-based incentives} (R\&D credits): 1,224 rows, 51 countries, 2019--2024
\end{itemize}

These forward-looking rates measure the tax burden on a hypothetical investment project, complementing the backward-looking CbCR ETRs. They are useful for analysing the incentive structure of tax systems but not directly applicable for estimating actual GMT exposure.

\subsection{Wage and earnings data}\label{app:wages}

Labour cost data are collected from three sources:

\begin{itemize}[noitemsep]
\item \textbf{OECD average annual wages}: 3,924 rows, 41 countries, 1990--2024. Average annual wages of full-time dependent employees in USD~PPP (2024 constant).
\item \textbf{OECD hourly earnings index}: 41,421 rows, 41 countries, 1955--2026. Index (2015~=~100); for trends rather than levels.
\item \textbf{ILO hourly earnings}: 227,317 rows, 105 countries, 1990--2025. Mean nominal hourly earnings in local currency and PPP~USD.
\item \textbf{ILO monthly earnings}: 529,234 rows, $\sim$140 countries, 1990--2025. Extends coverage to non-OECD economies.
\end{itemize}

OECD wage data is preferred for OECD countries; ILO data fills the gap for non-OECD members. These are used for payroll imputation (Assumption~A3).

\subsection{World Bank macroeconomic controls}\label{app:worldbank}

\begin{itemize}[noitemsep]
\item \textbf{GDP per capita} (\code{NY.GDP.PCAP.CD}): $\sim$2,660 rows, 2015--2024, $\sim$217 countries
\item \textbf{Population} (\code{SP.POP.TOTL}): $\sim$2,660 rows, 2015--2024
\item \textbf{Income group classifications}: 296 rows, cross-section (HIC/UMC/LMC/LIC)
\end{itemize}

\subsection{OECD BEPS implementation indicators}\label{app:beps}

\begin{itemize}[noitemsep]
\item \textbf{CbCR requirements} (Action~13): 493 rows, 119 countries
\item \textbf{CFC rules}: 1,901 rows, 146 countries, 2019/2024/2025
\item \textbf{Hybrid mismatch rules} (Action~2): 1,171 rows, 145 countries, 2025
\item \textbf{Interest limitation rules} (Action~4): 3,879 rows, 146 countries
\item \textbf{IP regimes} (Action~5): 2,305 rows, 64 countries, 2018--2025
\item \textbf{Mandatory disclosure rules} (Action~12): 711 rows, 145 countries, 2025
\end{itemize}

These provide policy indicator variables coding whether specific anti-avoidance laws are in place, their thresholds, and introduction dates.

\subsection{Pillar Two implementation status}\label{app:pillar2}

\textbf{PwC Pillar Two country tracker}: 1,190 rows (requirement $\times$ effective-date level). Key columns: Territory, Requirement (GIR, IIR, UTPR, QDMTT), Status (Enacted, Consultation, etc.), P2~Start date. Only countries with legislative steps appear. Cross-referenced against the OECD Central Record of Qualified Domestic Minimum Taxes (PDF).

\subsection{BEA US multinational enterprise data}\label{app:bea}

Four Excel workbooks on majority-owned foreign affiliates (MOFAs): selected data by country, employment, assets, and value added. Cover US-parented MNEs only. Useful for external validation and for the implementation regime analysis (Factor~4), given that US MNEs account for a disproportionate share of global profit shifting \citep{torslov2023}.

\subsection{Eurostat inward Foreign Affiliates Statistics}\label{app:fats}

\textbf{Combined rows:} 4,623,851 \quad \textbf{Combined size:} 442~MB \quad \textbf{Years:} 2008--2023

Eurostat inward FATS provide information on foreign-controlled enterprises operating inside European reporting economies. Key variables include number of enterprises, employment, turnover, value added, and personnel costs, disaggregated by NACE~Rev.~2 sector and controlling country. The data is mandatory from reference year 2021 onwards.

\textbf{Series break at 2021.} The legacy series \code{fats\_g1b\_08} (2008--2020) and the new series \code{fats\_activ} (2021--2023) differ in activity dimension (broad aggregate vs.\ 48 NACE categories), indicator codes (23 legacy vs.\ 19 revised), and control-country coverage (255 vs.\ 50 categories). Key concept changes: ``turnover'' becomes ``net turnover''; ``personnel costs'' becomes ``employee benefits expense''; the investment concept is narrowed to tangible non-current assets. Any analysis spanning both sides of the 2021 break must mark the break explicitly (see Assumption~A13a).

\subsection{OECD comparative tax revenue statistics}\label{app:cts}

\textbf{Rows:} 3,217,432 \quad \textbf{Size:} 342~MB \quad \textbf{Years:} 1990--2024 \quad \textbf{Countries:} 143

Harmonised time series of tax revenue by type and government level. The \code{T\_2000} series (\% of GDP) provides the baseline for assessing how GMT top-up revenues compare with existing CIT receipts.


%% ——————————————————————————————————————————————————————————————————
%%  WORKING DOCUMENT — NOT FOR SUBMISSION
%%  Set \workingdocfalse at the top to exclude this entire section.
%% ——————————————————————————————————————————————————————————————————
\ifworkingdoc

\newpage

\begin{center}
  \rule{\textwidth}{1pt}\\[0.8em]
  {\Large\bfseries Working Appendix}\\[0.4em]
  {\normalsize\itshape Not for submission - set \texttt{\textbackslash workingdocfalse} to exclude}\\[0.8em]
  \rule{\textwidth}{1pt}
\end{center}

\vspace{1em}

\renewcommand{\thesection}{W.\arabic{section}}
\setcounter{section}{0}
\setcounter{table}{0}
\setcounter{figure}{0}
\setcounter{equation}{0}

\begin{center}
\rule{0.8\textwidth}{1pt}\\[1em]
{\Large\bfseries Working Document}\\[0.5em]
{\normalsize Internal notes, redundant material, and drafting space.\\Remove entirely before final submission.}\\[0.5em]
\rule{0.8\textwidth}{1pt}
\end{center}

\addcontentsline{toc}{section}{Working Document (internal)}

%% ─────────────────────────────────────────────────────────────────────
%%  W.1: Introduction working notes
%% ─────────────────────────────────────────────────────────────────────
\section{Introduction --- working notes}\label{w:intro}

\subsection{Framing alternatives}

The thesis is framed as answering:

\begin{quote}
\textit{What are the revenue effects of the Global Minimum Tax for each EU Member State, and how do these estimates change when within-jurisdiction ETR heterogeneity, sectoral SBIE variation, and behavioural responses are incorporated?}
\end{quote}

The subquestions corresponding to the layers:
\begin{enumerate}
\item What is the mechanical top-up tax revenue for each EU Member State under the Pillar Two rule order (QDMTT/IIR)?
\item How much low-taxed profit exists within high-average-ETR EU jurisdictions, and how does incorporating this change the revenue estimate?
\item Which EU sectors are effectively shielded from the GMT by the SBIE, and how does the effective tax floor evolve over the transition period?
\item What is the plausible range of revenue effects when MNEs' behavioural responses to the higher effective tax rate are taken into account?
\end{enumerate}

\subsection{What the contribution is not}

The contribution does \emph{not} rest on a causal claim about the effects of GMT implementation. The data pre-dates full implementation (2022 CbCR vs.\ 2024 entry into force), and the behavioural response module uses imported parameters rather than identified variation. The paper is positioned as a measurement and simulation exercise, analogous to microsimulation studies in public economics.

%% ─────────────────────────────────────────────────────────────────────
%%  W.2: Literature review working notes
%% ─────────────────────────────────────────────────────────────────────
\section{Literature review --- working notes}\label{w:litreview}

\subsection{Constructed macro-level profit-shifting databases}

\citet{torslov2023} construct a global database of profits, wages and output corrected for profit shifting by combining three primary sources. First, they use national accounts data from the OECD's detailed sectoral accounts, supplemented by official statistics from the central banks and statistical institutes of non-OECD tax havens (Singapore, Hong Kong, Puerto Rico), to obtain corporate operating surplus, net interest payments and wages for all firms in each country. Second, they use foreign affiliates statistics (FATS) from Eurostat, the OECD and national statistical agencies. These census-type surveys record value-added, operating surplus and compensation of employees for all firms more than 50\% foreign-owned, tabulated by country of foreign parent. By combining FATS with national accounts, the authors compute profitability ratios (pre-tax profits per dollar of wages) separately for foreign and local firms in each jurisdiction. The excess profitability of foreign firms in tax havens relative to local firms is then attributed to profit shifting. Third, they use bilateral balance-of-payments data from the IMF and Eurostat to allocate shifted profits to source and parent countries, relying on bilateral service trade flows and intra-group interest payments to trace the geography of profit shifting. The resulting dataset, publicly available at \texttt{missingprofits.world}, covers the year 2015 and estimates that 36\% of multinational profits (\$616 billion) were shifted to tax havens globally.

\citet{barake2022} build on this data infrastructure to estimate the revenue potential of the GMT. They combine OECD CbCR aggregated statistics with the \citet{torslov2023} data to obtain profits, effective tax rates and payroll by headquarters-jurisdiction pair. For the substance-based income exclusion, tangible assets are drawn from CbCR where available, while the payroll component is approximated as the number of CbCR-reported employees multiplied by local mean earnings from ILO wage data. Revenue gains from non-haven jurisdictions are imputed using the ratio of non-haven to haven undertaxed profits, assuming a 20\% ETR in non-havens. These data choices enable a global revenue estimate of \EUR 179~billion (or \EUR 139--165~billion with carve-outs).

\subsection{Research gaps identified in literature overview}

\begin{enumerate}
\item \textbf{Reconciling macro and micro estimates.} The divergence between \citet{torslov2023} and \citet{blouin2023} remains unresolved. No study has yet attempted to reconcile these approaches using post-2020 data.

\item \textbf{Revenue estimates incorporating within-jurisdiction ETR heterogeneity and nonlinear behavioural responses.} No study has combined the within-jurisdiction ETR distribution data of \citet{hugger2026} with the nonlinear profit-shifting elasticities of \citet{dowd2017} and \citet{garcia-bernardo2024}.

\item \textbf{The SBIE as a production subsidy.} \citet{schjelderup2024} demonstrate theoretically that the SBIE functions as a production subsidy, but no empirical study has tested this prediction.

\item \textbf{Early empirical effects of GMT implementation.} The GMT was implemented in 2024/2025, creating a natural experiment of enormous scale. Despite this, no empirical study has yet estimated the initial effects. Timeline for early effects lies after summer 2026.

\item \textbf{Enforcement cooperation under the GMT.} The enforcement dilemma identified by \citet{jean2026} is testable in principle but data on MAP is scarce and unreliable.
\end{enumerate}

\subsection{Option B notes (not pursued but useful for later)}

\textbf{Research question:} Did the announcement of the GMT reduce the share of multinational profits reported in low-tax jurisdictions, and did EU-headquartered MNEs show a stronger anticipatory response than non-EU-headquartered MNEs?

This option carries the highest execution risk: the post-announcement window is narrow (only 2021--2022), COVID contaminates 2020--2021, and the effective sample size limits demanding subgroup analyses. Despite these risks, this option has the strongest causal identification. Even a null result would be informative. This remains an option for a later paper.

%% ─────────────────────────────────────────────────────────────────────
%%  W.3: Theoretical framework working notes
%% ─────────────────────────────────────────────────────────────────────
\section{Theoretical framework --- working notes}\label{w:theory}

\subsection{Double deduction derivation (extended)}

Under Pillar Two, the total tax on an MNE subsidiary in a low-tax jurisdiction ($\tau_{p,j} < T$) with positive excess profit is:

\begin{align}
\text{Tax} &= \tau_{p,j} \cdot \pi_{p,j} + T \cdot \pi_{p,j} - T \cdot \text{SBIE}_{p,j} - \tau_{p,j} \cdot \pi_{p,j} + \tau_{p,j} \cdot \text{SBIE}_{p,j} \nonumber \\
&= T \cdot \pi_{p,j} - (T - \tau_{p,j}) \cdot \text{SBIE}_{p,j}
\end{align}

For the numerical example from \citet{schjelderup2024}: $\pi = 1000$, $\tau = 5\%$, $\text{SBIE} = 650$ (at initial rates), $\delta^{\text{top}} = 10\%$, $\rho = 65\%$. At long-run rates with lower SBIE: effective floor rises.

\subsection{Key formulas from literature}

\begin{itemize}
\item Barake: $\text{Revenue}_p = \sum_j (0.15 - \text{ETR}) \times \max(0, \text{Profit} - \text{SBIE})$
\item Dowd quadratic SE: $\varepsilon(\tau) = -10.73 + 16.178(1-\tau)$ [statutory rates]
\item Garcia-Bernardo log SE: $\varepsilon(\tau) = \beta_3 + \beta_4/(t+\tau)$, $t \approx 0.002$
\item Schjelderup effective floor: $\tau^{P2} = T - (T-\tau) \times \text{SBIE}/\text{Profit}$
\item Devereux SBIE/profit ratio: 57\% average (year 1), 41\% (year 10)
\end{itemize}

%% ─────────────────────────────────────────────────────────────────────
%%  W.4: Methodology working notes
%% ─────────────────────────────────────────────────────────────────────
\section{Methodology --- working notes}\label{w:methodology}

\subsection{v1 layered architecture (superseded by v2 decomposition)}

The v1 methodology was structured as four cumulative layers, each building on the previous one:

\begin{table}[H]
\centering
\caption{Overview of v1 layered architecture (superseded)}
\label{tab:layers_v1}
\small
\begin{tabular}{@{}clp{5.5cm}p{4.5cm}@{}}
\toprule
Layer & Name & What it adds & Key assumption \\
\midrule
0 & Mechanical baseline & Updated Barake-style revenue estimate using 2022 CbCR & Jurisdiction-average ETRs are representative \\
1 & Within-jurisdiction heterogeneity & Replaces average ETRs with estimated ETR distributions & Uniform distribution between percentiles \\
2 & SBIE decomposition & Decomposes SBIE into sectoral components & FATS sectoral shares proxy CbCR composition \\
3 & Behavioural response & Simulates profit reallocation using imported semi-elasticities & Structural stability of parameters \\
\bottomrule
\end{tabular}
\end{table}

The v2 methodology restructures this as a $2^5$ factorial decomposition, which is superior because it makes all modelling choices simultaneously visible and exactly decomposable.

\subsection{Extended assumptions from v1}

\assumption{A0.1v1}{Tax accrued (\code{TAX\_ACCRUED}) is the appropriate numerator for the CbCR ETR. This follows \citet{hugger2026}, who argue that accrued taxes reflect the current-year charge on the year's profits, unlike taxes paid (cash basis) which include timing differences. The GloBE ETR under Pillar Two uses ``covered taxes'' which differ from both accrued and paid taxes in CbCR, but CbCR accrued taxes are the closest available proxy. As \citet{hanlon2023} emphasises, GloBE income is neither taxable income nor financial accounting income; it requires at least 22 adjustments from the financial accounting starting point, including a deferred tax adjustment computed at the lower of 15\% or the statutory rate. The CbCR-based ETR is therefore an approximation of the GloBE ETR, and the direction of bias depends on the local tax system: jurisdictions with generous accelerated depreciation may have lower CbCR ETRs than GloBE ETRs (since the deferred tax adjustment partially reverses timing differences), while jurisdictions relying on qualified refundable tax credits---which under GloBE rules are treated as income rather than as a tax reduction \citep{ferreiraliotti2022}---may have higher CbCR ETRs than GloBE ETRs.}

\assumption{A0.7v1}{Where \code{PROFIT\_ADJ} is available, it is used as the profit measure for the ETR calculation. Where it is not available, the unadjusted \code{PROFIT} is used with a correction factor derived from the ratio of adjusted to unadjusted profits across jurisdictions where both are observed. This follows the approach of \citet{hugger2026}. The importance of this correction cannot be overstated. \citet{blouin2023} demonstrate that equity method accounting in aggregate data causes systematic double-counting of MNE profits: income flows up through ownership chains, and each holding entity records the subsidiary's profit as its own. In the US BEA data, correcting for equity income reduced \citeauthor{clausing2016}'s (\citeyear{clausing2016}) US revenue loss estimate from \$77--111 billion to \$11 billion---a factor of ten. The Netherlands, conventionally classified as a top US tax haven with a 4.4\% ETR, has a corrected ETR of 25\% once equity income is excluded. For CbCR data specifically, \citet{garcia-bernardo2022} find that CbCR overstated total US corporate profits by 47--72\% in the years 2017--2019, and \citet{hugger2026} apply a median correction of 32.4\% across jurisdictions.}

\subsection{Additional assumptions from v2 not in main text}

\assumption{A2a}{(CIT-exempt entities). CbCR aggregates include CIT-exempt entities (pension funds, sovereign wealth funds, universities) unless the reporting jurisdiction specifies otherwise \citep{oecd2025a}. Under GloBE rules, these entities are excluded from Pillar Two scope (Article~1.5). Their inclusion depresses the measured ETR, introducing a downward bias.}

\assumption{A4a}{(Qualified refundable tax credits). CbCR does not distinguish between qualified refundable tax credits (QRTCs) and non-qualified credits. Under GloBE rules, QRTCs are treated as income rather than as a reduction of covered taxes, raising the GloBE ETR relative to the CbCR ETR \citep{ferreiraliotti2022}. Jurisdictions relying heavily on QRTCs---notably Canada and France---may appear low-taxed in CbCR data but face higher GloBE ETRs.}

\assumption{A4b}{(Minority ownership). The IIR applies in proportion to the UPE's ownership share. CbCR data report 100\% of each entity's financial variables regardless of ownership fraction. \citet{barake2022} note that approximately 8.7\% of MNE profits are minority-owned; treating all profit as fully owned overstates the IIR tax base by up to 8--9\%.}

\assumption{A4c}{(Loss carry-forward). GloBE rules allow net losses to reduce GloBE income in subsequent years. CbCR panel~AI excludes loss-making entities but does not track loss carry-forwards. The methodology applies a static year-by-year computation. The bias is bounded by loss-adjustment factors: positive profits are reduced by 11.8\% on average (18.1\% for investment hubs).}

\assumption{A4d}{(Debt-financing bias). The SBIE-to-profit ratio is mechanically higher for subsidiaries with aggressive debt structures, because interest deductions reduce pre-tax profit while leaving payroll and tangible assets unchanged. \citet{devereux2023} confirm this empirically. The methodology does not adjust for this interaction.}

\subsection{Feasibility taxonomy}

\begin{table}[H]
\centering
\caption{Feasibility taxonomy}
\label{tab:feasibility}
\small
\begin{tabular}{@{}p{4cm}ccp{5cm}@{}}
\toprule
Component & Data ready? & Assumption burden & Execution risk \\
\midrule
\multicolumn{4}{l}{\textit{Baseline computation}} \\
ETR calculation & Yes & Low & Low \\
SBIE aggregate & Yes & Low--Medium & Low \\
Revenue allocation & Yes & Low & Low \\
Double-counting correction & Partial & Medium & Medium \\
\addlinespace
\multicolumn{4}{l}{\textit{ETR distribution}} \\
ETR percentiles & Yes & Medium & Low--Medium \\
Piecewise-uniform fitting & Yes & Medium & Low \\
Cross-validation with Table IV & Yes & Low & Low \\
\addlinespace
\multicolumn{4}{l}{\textit{SBIE sectoral decomposition}} \\
FATS sectoral payroll shares & Yes & Medium & Medium \\
FATS sectoral tangible assets & No (imputed) & High & High \\
Effective tax floor by sector & Derived & Medium & Medium \\
\addlinespace
\multicolumn{4}{l}{\textit{Behavioural response}} \\
Linear semi-elasticity & Imported & Low & Low \\
Quadratic semi-elasticity & Imported & Medium & Medium \\
Logarithmic semi-elasticity & Imported & High & High \\
\bottomrule
\end{tabular}
\end{table}

\subsection{What to prune}

\begin{itemize}
\item \textbf{Causal identification of GMT announcement effects.} CbCR covers only up to 2022 (one year post-announcement). Separate paper for when 2023--2024 data becomes available.
\item \textbf{Developing country distributional analysis.} Data coverage outside the EU is weaker.
\item \textbf{Enforcement cooperation effects.} Data on MAP is scarce and unreliable.
\item \textbf{Full Orbis firm-level analysis as core.} Use Orbis only for validation.
\item \textbf{Re-estimation of semi-elasticity parameters.} Not credible from aggregate CbCR data.
\end{itemize}

%% ─────────────────────────────────────────────────────────────────────
%%  W.5: What the methodology does not model (extended)
%% ─────────────────────────────────────────────────────────────────────
\section{What the methodology does not model --- extended notes}\label{w:limitations}

\paragraph{Real economic responses.} The semi-elasticity specifications model reallocation of \emph{reported} profits (paper shifting), but the GMT may also trigger reallocation of real economic activity. \citet{suarezserrato2018} find that the repeal of US Section~936 reduced global investment by 10\% and US employment by 6.7\% for exposed firms, with a capital semi-elasticity of 1.77 and a local labour market multiplier of 3.87. \citet{bilicka2022} find that the UK Worldwide Debt Cap reduced total assets by 30.6\% and employment by 11.3\% among treated foreign MNEs. If the GMT triggers analogous responses, the SBIE itself would shrink, amplifying the top-up tax in a second-round effect not captured here.

\paragraph{Government strategic responses.} \citet{janeba2023} show that in a Nash equilibrium with distortionary tax competition, non-haven governments do \emph{not} reduce tax rates after GMT implementation---they may even raise them, because the GMT reduces profit shifting and thereby increases the value of attracting real FDI, producing revenue gains beyond the mechanical effect. However, if governments compete via firm-specific lump-sum subsidies rather than tax rates, the GMT leaves net revenue unchanged while haven revenue increases. The introduction of QDMTTs in Ireland, the Netherlands, and Luxembourg is itself a strategic response. \citet{barake2022} show that under a pure QDMTT regime, US revenue from Pillar Two would collapse from approximately EUR~54 billion (under IIR) to EUR~3 billion.

\paragraph{Enforcement dilemma.} \citet{jean2026} model the interaction between the GMT and endogenous enforcement effort. There exists a threshold beyond which the low-tax country prefers to defect from enforcement cooperation. Since the EU relies on automatic information exchange and mutual administrative assistance, the GMT may erode the cooperative infrastructure that currently constrains profit shifting.

\paragraph{Welfare analysis.} Revenue estimates do not map directly to welfare. \citet{johannesen2022} demonstrates theoretical ambiguity: the net welfare effect is positive only when the profit-shifting semi-elasticity is sufficiently high. The consensus estimate of 0.8 falls well below the threshold of approximately~4 required for the GMT to be unambiguously welfare-improving, though the quadratic specification of \citet{dowd2017} yields semi-elasticities above~4 for the lowest-taxed jurisdictions. \citet{hebous2023} show that a carefully designed minimum tax can be Pareto-improving, but only under specific conditions.

\paragraph{Profit-shifting channels.} The methodology does not distinguish between profit-shifting mechanisms. \citet{torslov2023} decompose shifted profits into transfer pricing and IP location ($\sim$85\%) versus intra-group debt ($\sim$15\%). The GMT affects these channels differently.

\paragraph{Developing country effects.} When an EU parent collects IIR top-up tax on a low-taxed affiliate in a developing country without a QDMTT, the revenue accrues to the EU Member State rather than the source country. \citet{riet2025} show that developing countries that offer tax incentives lose both the investment incentive and the tax revenue, unless they implement a QDMTT.

%% ─────────────────────────────────────────────────────────────────────
%%  W.6: Data annex working notes
%% ─────────────────────────────────────────────────────────────────────
\section{Data annex --- working notes}\label{w:data}

\subsection{OECD CbCR data structure}

Under the OECD/G20 BEPS Action~13, large MNE groups with consolidated revenue exceeding EUR~750~million are required to file Country-by-Country Reports detailing their economic activity and tax payments in each jurisdiction. The OECD aggregates and anonymises these reports into jurisdiction-level statistics published in the \emph{OECD Corporate Tax Statistics} database. Every dataset is identified by a \emph{dataflow} code of the form \code{OECD.CTP.TPS,DSD\_CBCR@DF\_CBCR[I--VI],1.1}. The tables cover fiscal years 2016--2022.

\subsection{Common column structure}

All six files share fifteen columns: \code{DATAFLOW}, \code{REF\_AREA} (parent jurisdiction), \code{COUNTERPART\_AREA} (affiliate jurisdiction), \code{FREQ} (always annual), \code{MEASURE}, \code{UNIT\_MEASURE}, \code{STATISTICAL\_OPERATION}, \code{AGGREGATION\_TYPE}, \code{PROFIT\_GROUPING}, \code{EFFECTIVE\_TAX\_RATE}, \code{MNE\_SIZE}, \code{TIME\_PERIOD}, \code{OBS\_VALUE}, \code{DECIMALS}, \code{OBS\_STATUS}.

\subsection{FATS series break}

Eurostat inward FATS underwent a structural break in 2021. The pre-2021 dataset (\code{fats\_g1b\_08}) uses 23 legacy indicator codes, 255 control-country categories, and a broad activity aggregate (\code{B-N\_S95\_X\_K}). The post-2021 dataset (\code{fats\_activ}) uses 19 revised indicator codes, 50 control-country categories, and a different aggregate (\code{B-S\_X\_O\_S94}). Key concept changes include ``turnover'' becoming ``net turnover'' and ``personnel costs'' becoming ``employee benefits expense.'' The sectoral analysis uses the post-2021 structure for all cross-sectional results.

\subsection{Switzerland anomaly}

Switzerland has anomalous values in Table~IV: profit values approximately 8,600$\times$ larger than Table~I in 2017, 2020, and 2022. Exclude CHE from Table~IV cross-validation.

%% ─────────────────────────────────────────────────────────────────────
%%  W.7: Results working notes
%% ─────────────────────────────────────────────────────────────────────
\section{Results --- working notes}\label{w:results}

[Space for preliminary calculations, draft tables, and intermediate results.]

%% ─────────────────────────────────────────────────────────────────────
%%  W.8: Discussion working notes
%% ─────────────────────────────────────────────────────────────────────
\section{Discussion --- working notes}\label{w:discussion}

\subsection{Target journals and positioning}

The methodology is designed to be compatible with publication in applied public economics or tax policy journals. The side-by-side architecture and feasibility taxonomy are suited to journals that value methodological transparency and policy relevance (e.g.\ \emph{International Tax and Public Finance}, \emph{Fiscal Studies}, \emph{Intertax}). The EU focus and sectoral SBIE analysis also connect to the European policy debate around the Pillar Two Directive implementation.

%% ─────────────────────────────────────────────────────────────────────
%%  W.7a: Thought process for March 2026 revisions
%% ─────────────────────────────────────────────────────────────────────
\section{Thought process: March 2026 revisions}\label{w:thought_process}

This section documents the reasoning behind the major structural changes made in March 2026 to the theoretical framework and methodology.

\subsection{Why the revenue pipeline model}

The original Section~3 had three disconnected parts: the Schjelderup effective floor derivation, the sectoral decomposition, and a list of four hypotheses. A referee or examiner would ask: ``What is the unifying model that generates these hypotheses?'' The literature offers several candidates:

\begin{itemize}[noitemsep]
\item \textbf{Full general equilibrium} (Janeba \& Schjelderup 2023): too complex for a thesis that is fundamentally a microsimulation exercise. Would require calibrating strategic responses of governments, which the data cannot support.
\item \textbf{Welfare-theoretic model} (Johannesen 2022): focuses on optimal policy rather than revenue forecasting. The semi-elasticity threshold result ($\varepsilon > 4$) is incorporated as a welfare diagnostic but does not structure the revenue computation.
\item \textbf{Accounting decomposition}: the simplest option that is sufficient for the purpose. The revenue pipeline (equation~\eqref{eq:pipeline}) is an accounting identity that decomposes revenue into three multiplicative filters. It does not claim behavioural content; it structures the measurement exercise.
\end{itemize}

The pipeline was chosen because it (a)~maps directly onto the factorial decomposition, making the theory-to-empirics connection explicit; (b)~is general enough to accommodate future extensions (e.g., adding a Stage~4 for government strategic responses); and (c)~is honest about the nature of the exercise: this is a measurement and simulation study, not a structural estimation.

\subsection{Why the descriptive regression}

The original methodology was 100\% microsimulation---no regression analysis. This is a problem for three reasons:

\begin{enumerate}[noitemsep]
\item \textbf{Publication convention}: applied public economics journals expect at least some regression-based evidence. A pure simulation exercise reads as a technical report rather than an academic paper.
\item \textbf{Testing H2}: the SBIE sectoral asymmetry hypothesis is most naturally tested in a regression framework, where the manufacturing share coefficient provides a single summary statistic with a standard error.
\item \textbf{Comparability}: Hugger et al.\ (2025) run a cross-sectional regression predicting average ETR. Running an analogous regression predicting the effective floor shows how the thesis's novel dependent variable relates to established work.
\end{enumerate}

The regression is explicitly descriptive. The effective floor is a model output, not an observed outcome. The regression summarises cross-sectional correlates, not causal effects. This is stated clearly in the text and footnote.

\subsection{Why H5 was added}

The Table~IV cross-validation was originally presented as a methodological exercise without a prior prediction. This is weak: a contribution should have a theoretical anchor. H5 predicts that conduit jurisdictions will show $\lambda > 1$ because their bimodal ETR distributions are compressed by the piecewise-uniform assumption. The prediction is falsifiable: if $\lambda \approx 1$ for conduits, the piecewise-uniform assumption is adequate and H5 is rejected.

\subsection{Why H3 was reframed as corollary of H2}

The original thesis had four hypotheses for three research questions. H3 (Ireland--Netherlands divergence) is logically a special case of H2 (SBIE sectoral asymmetry). Presenting it as an independent hypothesis inflates the hypothesis count and obscures the logical structure. Reframing it as ``an application of H2'' preserves the analysis but clarifies that the Ireland--Netherlands comparison is an illustrative test, not an independent prediction.

\subsection{Assessment of what was pruned}

\begin{itemize}[noitemsep]
\item \textbf{Causal identification}: Not feasible with 2016--2022 data (only one post-announcement year). Saved for future paper with 2023--2024 data.
\item \textbf{Welfare computation}: Deferred to a footnote (Johannesen threshold). A full welfare analysis requires calibrating redistributive preferences, which is beyond the scope.
\item \textbf{Enforcement cooperation} (Jean \& Nishimura 2026): Mentioned in the literature review and limitations, but not modelled. Data on mutual administrative assistance is unreliable.
\item \textbf{Pillar One interaction}: Noted in limitations. Modelling Amount~A reallocation would require a separate data infrastructure.
\item \textbf{Re-estimation of semi-elasticities}: CbCR aggregate data lacks the within-firm variation required for identification. Importing parameters is the honest approach.
\end{itemize}

\subsection{Literature cross-check: which papers informed which additions}

\begin{itemize}[noitemsep]
\item \textbf{Revenue pipeline}: Novel formulation. Components drawn from Hugger (Stage~1), Schjelderup/Devereux (Stage~2), Dowd/Garcia-Bernardo (Stage~3). No prior study combines these multiplicatively.
\item \textbf{Bootstrap CIs}: Standard in microsimulation literature (e.g., Clausing 2016 uses sensitivity ranges but not bootstrap). The bootstrap on bilateral pairs follows the resampling approach in the microeconometrics literature.
\item \textbf{Shapley decomposition}: Follows Huettner \& Sunder (2012) for decomposing explained variance. Applied here to revenue variation rather than $R^2$.
\item \textbf{Descriptive regression}: Specification parallels Hugger (2025) Table~4 (average ETR regression), with novel dependent variable (effective floor) and manufacturing share as key predictor.
\item \textbf{Table~IV activity composition}: Uses CbCR variables documented in OECD (2025) metadata. The cross-check against FATS is novel.
\item \textbf{Time-series dimension}: Straightforward extension of the cross-sectional analysis to the 2016--2022 panel. No prior study has tracked the GMT base over this full period.
\item \textbf{Welfare threshold}: Direct application of Johannesen (2022) Proposition~3 to the Dowd quadratic semi-elasticity. The 7.7\% threshold is computed from $\varepsilon_Q(\tau) = 4 \Rightarrow \tau = 1 - 14.73/16.178 \approx 0.077$.
\end{itemize}

\section{TODO list}\label{w:todos}

This section tracks all items requiring additional research, writing, or verification before submission.

\subsection{Writing tasks}

\begin{enumerate}
\item \textbf{Abstract} --- Write the abstract (currently placeholder). Should summarise all three research questions, the $2^5$ factorial decomposition approach, and headline results.
\item \textbf{Introduction contribution paragraph} --- The \texttt{\textbackslash todo} notes flag that the contribution needs to be stated more clearly in the introduction. Merge the Contribution subsection (\ref{sec:contribution}) content into a tighter paragraph.
\item \textbf{Results section} --- All subsections are placeholders. Priority order: (i)~descriptive statistics, (ii)~baseline revenue estimates (Approach~A vs.~B), (iii)~cross-validation, (iv)~decomposition results, (v)~SBIE sectoral decomposition, (vi)~behavioural sensitivity, (vii)~robustness checks.
\item \textbf{Discussion section} --- Position findings relative to \citet{barake2022}, \citet{hugger2026}, and OECD projections. Discuss implications for the SBIE design debate. Discuss QDMTT vs.\ IIR allocation consequences.
\item \textbf{Conclusion} --- Formulate answer to three research questions. Summarise theoretical and practical relevance. Nuance conclusions. Suggestions for future research.
\item \textbf{Limitations section} --- Expand the placeholder to cover: CbCR-to-GloBE approximation, no real economic responses, no government strategic responses, no enforcement dynamics, no Pillar One interaction, static year-by-year computation without loss carry-forward, no welfare analysis.
\end{enumerate}

\subsection{Data and computation tasks}

\begin{enumerate}
\item \textbf{Run baseline revenue computation} (Approach~A) for 2022 CbCR data. Compare against \citet{barake2022} EUR~179~billion estimate.
\item \textbf{Implement ETR distribution fitting} (Approach~B) using Table~I percentiles. Compute $s_{j,t}^{<T}$ for all jurisdictions.
\item \textbf{Cross-validate} Table~I distribution-implied low-taxed profit against Table~IV \code{RATE\_0-15} direct measurement. Document ratios by jurisdiction.
\item \textbf{Compute $2^5$ factorial decomposition.} All 32 scenarios. Extract marginal contributions, interactions, and factor shares.
\item \textbf{SBIE sectoral decomposition.} Merge FATS sectoral payroll shares with CbCR aggregates for EU-27. Compute sector-level SBIE-to-profit ratios.
\item \textbf{Tangible asset imputation.} Source OECD STAN capital intensity ratios by sector. This is the weakest link (Assumption~A14)---consider alternative imputation strategies.
\item \textbf{Effective tax floor matrices.} Compute for 27~EU $\times$ 15--20~NACE sectors $\times$ 10~transition years.
\item \textbf{Ireland--Netherlands comparison.} Extract sectoral composition from FATS. Test H3 (equation~\ref{eq:irl_nld_hypothesis}).
\item \textbf{Behavioural sensitivity.} Implement all three semi-elasticity specifications. Compute profit reallocation.
\item \textbf{Revenue uncertainty map.} Produce choropleth showing dominant factor and dominance ratio by jurisdiction.
\item \textbf{Double-counting correction.} Compute \code{PROFIT\_ADJ}/\code{PROFIT} ratios for all jurisdictions where both available. Test sensitivity of results.
\item \textbf{Loss adjustment.} Implement \citet{hugger2026} loss-adjustment factors (11.8\% average, 18.1\% for investment hubs).
\end{enumerate}

\subsection{Research and verification tasks}

\begin{enumerate}
\item \textbf{Verify citation keys.} The \citet{jean2026} reference appears as ``jean2025'' in the literature analyses folder---confirm correct publication year. Similarly verify \citet{hugger2026} vs.\ ``hugger2025'' (OECD working paper vs.\ published version).
\item \textbf{Verify Dowd formula application.} The quadratic semi-elasticity is estimated on \emph{statutory} rates but applied to \emph{effective} rates (Assumption~A16a). Quantify the bias for key jurisdictions (Ireland, Netherlands, Luxembourg).
\item \textbf{Check OECD Investment Tax Incentives Database} for jurisdictions requiring asymmetric shift in Step~3 of Approach~B.
\item \textbf{QDMTT implementation status.} Update PwC tracker data to current status. Verify all EU-27 have enacted QDMTTs.
\item \textbf{Side-by-Side (SbS) package details.} Verify January 2026 OECD/G20 agreement terms for US-parented MNE groups. Confirm IIR/UTPR exemption scope.
\item \textbf{FATS-CbCR proxy validation.} Compare FATS total employment against CbCR \code{EMPLOYEES} for overlapping EU jurisdictions. Document any systematic discrepancy.
\item \textbf{Table~IV activity composition analysis.} Extract \code{ACT\_HOLDING}, \code{ACT\_IP}, \code{ACT\_MAN} shares for the \code{RATE\_0-15} bracket to validate FATS-based sectoral decomposition.
\item \textbf{Additional literature to consider incorporating:}
  \begin{itemize}[noitemsep]
  \item \citet{beer2020} meta-analysis on semi-elasticity stability (already referenced in A15)
  \item \citet{schanz2023} on GMT compliance costs and safe harbours (referenced in A18)
  \item \citeauthor{boadway2026} on conceptual parallels with global minimum wealth tax
  \item \citeauthor{akcigit2022} on R\&D policy interactions
  \item Brun (2025) on GMT revenue simulation
  \end{itemize}
\item \textbf{Presentation of implementation map.} Create a figure or table showing which jurisdictions have implemented IIR, QDMTT, and UTPR, and when.
\item \textbf{Framing for non-fiscal readers.} The introduction \texttt{\textbackslash todo} notes that the arm's-length principle reference may not be accessible to non-fiscal readers. Consider reframing the opening in more general economic policy terms.
\end{enumerate}

\subsection{Figures and tables to produce}

\begin{enumerate}
\item Revenue uncertainty map (choropleth, Section~\ref{sec:uncertainty_map})
\item Factor decomposition bar chart (EU-27 and worldwide aggregates)
\item Effective tax floor heatmap (EU-27 $\times$ NACE sectors)
\item SBIE transition path (2024--2033) for selected jurisdictions
\item Ireland vs.\ Netherlands sectoral comparison
\item Table~I vs.\ Table~IV cross-validation scatter plot
\item Scenario matrix (Table~\ref{tab:scenario_matrix}) with computed values
\item EU-27 Member State revenue table (QDMTT + IIR, net fiscal effect)
\item ETR distribution plots for selected jurisdictions (P5--P95)
\item Time series of low-taxed profit share (2016--2022)
\end{enumerate}

\citet{bilicka2023} extend this by examining tax-avoidance networks through the lens of intellectual property, showing that a disproportionate share of patents (particularly high-value patents) are registered in tax havens, consistent with the strategic location of IP being a primary channel for profit shifting. Five percent of European patent applications end up in tax havens, and 20\% of global cross-border patent transactions involve tax haven buyers. MNEs that meet the size threshold requirements for the proposed GMT are particularly active in developing patents: they constitute 2.6\% of affiliates but are responsible for 42\% of all patent applications and 45\% of tax haven ones. Patenting is concentrated in the top quantile of the turnover distribution suggesting that IP-intensive activities, which generate little SBIE, are particularly exposed to the GMT.

\subsection{Welfare analysis (for discussion section)}

\citet{vantriet2025} assess the welfare implications of the GMT in a multicountry macroeconomic model. They simulate the implementation of a GMT of 15\% by all countries, finding that when additional CIT revenues are redistributed as direct transfers to households, the global welfare impact is slightly positive but mixed across countries. When additional revenues are redistributed back to firms as lower CIT rates (provided the rate remains at or above the GMT rate), positive welfare outcomes are more widely experienced. They find that a 16\% GMT rate yields the highest level of global welfare in their model.

\fi  % end \ifworkingdoc

\end{document}
